\chapter{Repository Administration}\label{svn.reposadmin}

The Subversion repository is the central storehouse of all your
versioned data. As such, it becomes an obvious candidate for all the
love and attention an administrator can offer. While the repository is
generally a low-maintenance item, it is important to understand how to
properly configure and care for it so that potential problems are
avoided, and so actual problems are safely resolved.

In this chapter, we\textquotesingle ll discuss how to create and
configure a Subversion repository. We\textquotesingle ll also talk about
repository maintenance, providing examples of how and when to use
various related tools provided with Subversion. We\textquotesingle ll
address some common questions and mistakes and give some suggestions on
how to arrange the data in the repository.

If you plan to access a Subversion repository only in the role of a user
whose data is under version control (i.e., via a Subversion client), you
can skip this chapter altogether. However, if you are, or wish to
become, a Subversion repository administrator,\footnote{This may sound
  really prestigious and lofty, but we\textquotesingle re just talking
  about anyone who is interested in that mysterious realm beyond the
  working copy where everyone\textquotesingle s data hangs out.} this
chapter is for you.

\section{The Subversion Repository,
Defined}\label{svn.reposadmin.basics}

Before jumping into the broader topic of repository administration,
let\textquotesingle s further define what a repository is. How does it
look? How does it feel? Does it take its tea hot or iced, sweetened, and
with lemon? As an administrator, you\textquotesingle ll be expected to
understand the composition of a repository both from a literal, OS-level
perspective---how a repository looks and acts with respect to
non-Subversion tools---and from a logical perspective---dealing with how
data is represented \emph{inside} the repository.

Seen through the eyes of a typical file browser application (such as
Windows Explorer) or command-line based filesystem navigation tools, the
Subversion repository is just another directory full of stuff. There are
some subdirectories with human-readable configuration files in them,
some subdirectories with some not-so-human-readable data files, and so
on. As in other areas of the Subversion design, modularity is given high
regard, and hierarchical organization is preferred to cluttered chaos.
So a shallow glance into a typical repository from a nuts-and-bolts
perspective is sufficient to reveal the basic components of the
repository:

\begin{verbatim}
$ ls repos
conf/  db/  format  hooks/  locks/  README.txt
\end{verbatim}

Here\textquotesingle s a quick fly-by overview of what exactly
you\textquotesingle re seeing in this directory listing.
(Don\textquotesingle t get bogged down in the terminology---detailed
coverage of these components exists elsewhere in this and other
chapters.)

\begin{description}
\item[conf/]
This directory is a container for configuration files.
\item[db/]
This directory contains the data store for all of your versioned
data.\footnote{Strictly speaking, Subversion doesn\textquotesingle t
  dictate that the versioned data live here, and there are known (albeit
  proprietary) alternative backend storage implementations which do not,
  in fact, store data in this directory.}
\item[format]
This file describes the repository\textquotesingle s internal
organizational scheme. (As it turns out, the \texttt{db/} subdirectory
sometimes also contains a \texttt{format} file which describes only the
contents of that subdirectory and which is not to be confused with this
file.)
\item[hooks/]
This directory contains hook script templates and hook scripts, if any
have been installed.
\item[locks/]
Subversion uses this directory to house repository lock files, used for
managing concurrent access to the repository.
\item[README.txt]
This is a brief text file containing merely a notice to readers that the
directory they are looking in is a Subversion repository.
\end{description}

Prior to Subversion 1.5, the on-disk repository structure also always
contained a \texttt{dav} subdirectory. \texttt{mod\_dav\_svn} used this
directory to store information about WebDAV activities---mappings of
high-level WebDAV protocol concepts to Subversion commit transactions.
Subversion 1.5 changed that behavior, moving ownership of the activities
directory, and the ability to configure its location, into
\texttt{mod\_dav\_svn} itself. Now, new repositories will not
necessarily have a \texttt{dav} subdirectory unless
\texttt{mod\_dav\_svn} is in use and hasn\textquotesingle t been
configured to store its activities database elsewhere. See
\hyperref[svn.serverconfig.httpd.ref.mod_dav_svn]{mod\_dav\_svn
configuration directives} for more information.

Of course, when accessed via the Subversion libraries, this otherwise
unremarkable collection of files and directories suddenly becomes an
implementation of a virtual, versioned filesystem, complete with
customizable event triggers. This filesystem has its own notions of
directories and files, very similar to the notions of such things held
by real filesystems (such as NTFS, FAT32, ext3, etc.). But this is a
special filesystem---it hangs these directories and files from
revisions, keeping all the changes you\textquotesingle ve ever made to
them safely stored and forever accessible. This is where the entirety of
your versioned data lives.

\protect\phantomsection\label{svn.reposadmin.basics.backends}
Speaking of Filesystems\ldots{}

When the initial design phase of Subversion was in progress, the
developers decided to use Berkeley DB (BDB) as the storage mechanism
behind the virtual versioned filesystem implementation. Berkeley DB was
a logical choice for a variety of reasons, including its open source
license, transaction support, reliability, performance, API simplicity,
thread safety, support for cursors, and so on.

In the years since, the newer FSFS\footnote{While it is often pronounced
  ``fuzz-fuzz,'' per Jack Repenning\textquotesingle s rendition, this
  book assumes that the reader is thinking ``eff-ess-eff-ess.''} backend
was introduced. This so-called ``filesystem filesystem'' was a versioned
filesystem implemented not within an opaque database container, but
instead as a larger collection of more transparent files stored in the
OS\textquotesingle s filesystem. FSFS enjoyed continual development and
improvement, and eventually earned the right to be the default
Subversion backend. But improvements to that backend kept coming, and
ultimately the FSFS storage layer surpassed the Berkeley DB one in
nearly every meaningful metric, from performance to scalability to
reliability and beyond.

These days, it is generally assumed that if you are using the open
source Subversion product, you are using the FSFS backend for your
repositories. In fact, beginning with Subversion 1.8, the Berkeley DB
Subversion repository filesystem backend has been officially deprecated.
Subversion repositories which still use this storage layer option will
continue to function with newer Subversion 1.x releases, but no further
development---including feature introduction or expansion---is planned
for the Berkeley DB backend. Subversion effectively offers a single
viable repository storage layer option. FSFS won.

This book assumes throughout what the rest of the world does: that FSFS
is \emph{the} Subversion storage backend implementation. Berkeley DB is
mentioned only for historical context; administrators of legacy
BDB-backed repositories should consult older versions of this
documentation for information about maintaining them.

\section{Strategies for Repository
Deployment}\label{svn.reposadmin.planning}

Due largely to the simplicity of the overall design of the Subversion
repository and the technologies on which it relies, creating and
configuring a repository are fairly straightforward tasks. There are a
few preliminary decisions you\textquotesingle ll want to make, but the
actual work involved in any given setup of a Subversion repository is
pretty basic, tending toward mindless repetition if you find yourself
setting up multiples of these things.

Some things you\textquotesingle ll want to consider beforehand, though,
are:

\begin{itemize}
\item
  What data do you expect to live in your repository (or repositories),
  and how will that data be organized?
\item
  Where will your repository live, and how will it be accessed?
\item
  What types of access control do you need?
\end{itemize}

In this section, we\textquotesingle ll try to help you answer those
questions.

\subsection{Planning Your Repository
Organization}\label{svn.reposadmin.projects.chooselayout}

While Subversion allows you to move around versioned files and
directories without any loss of information, and even provides ways of
moving whole sets of versioned history from one repository to another,
doing so can greatly disrupt the workflow of those who access the
repository often and come to expect things to be at certain locations.
So before creating a new repository, try to peer into the future a bit;
plan ahead before placing your data under version control. By
conscientiously ``laying out'' your repository or repositories and their
versioned contents ahead of time, you can prevent many future headaches.

Let\textquotesingle s assume that as repository administrator, you will
be responsible for supporting the version control system for several
projects. Your first decision is whether to use a single repository for
multiple projects, or to give each project its own repository, or some
compromise of these two.

There are benefits to using a single repository for multiple projects,
most obviously the lack of duplicated maintenance. A single repository
means that there is one set of hook programs, one thing to routinely
back up, one thing to dump and load if Subversion releases an
incompatible new version, and so on. Also, you can move data between
projects easily, without losing any historical versioning information.

The downside of using a single repository is that different projects may
have different requirements in terms of the repository event triggers,
such as needing to send commit notification emails to different mailing
lists, or having different definitions about what does and does not
constitute a legitimate commit. These aren\textquotesingle t
insurmountable problems, of course---it just means that all of your hook
scripts have to be sensitive to the layout of your repository rather
than assuming that the whole repository is associated with a single
group of people. Also, remember that Subversion uses repository-global
revision numbers. While those numbers don\textquotesingle t have any
particular magical powers, some folks still don\textquotesingle t like
the fact that even though no changes have been made to their project
lately, the youngest revision number for the repository keeps climbing
because other projects are actively adding new revisions.\footnote{Whether
  founded in ignorance or in poorly considered concepts about how to
  derive legitimate software development metrics, global revision
  numbers are a silly thing to fear, and \emph{not} the kind of thing
  you should weigh when deciding how to arrange your projects and
  repositories.}

A middle-ground approach can be taken, too. For example, projects can be
grouped by how well they relate to each other. You might have a few
repositories with a handful of projects in each repository. That way,
projects that are likely to want to share data can do so easily, and as
new revisions are added to the repository, at least the developers know
that those new revisions are at least remotely related to everyone who
uses that repository.

After deciding how to organize your projects with respect to
repositories, you\textquotesingle ll probably want to think about
directory hierarchies within the repositories themselves. Because
Subversion uses regular directory copies for branching and tagging (see
\hyperref[svn.branchmerge]{Branching and Merging}), the Subversion
community recommends that you choose a repository location for each
project root---the ``topmost'' directory that contains data related to
that project---and then create three subdirectories beneath that root:
\texttt{trunk}, meaning the directory under which the main project
development occurs; \texttt{branches}, which is a directory in which to
create various named branches of the main development line; and
\texttt{tags}, which is a collection of tree snapshots that are created,
and perhaps destroyed, but never changed.\footnote{The \texttt{trunk},
  \texttt{tags}, and \texttt{branches} trio is sometimes referred to as
  ``the TTB directories.''}

For example, your repository might look like this:

\begin{verbatim}
/
   calc/
      trunk/
      tags/
      branches/
   calendar/
      trunk/
      tags/
      branches/
   spreadsheet/
      trunk/
      tags/
      branches/
   …
\end{verbatim}

Note that it doesn\textquotesingle t matter where in your repository
each project root is. If you have only one project per repository, the
logical place to put each project root is at the root of that
project\textquotesingle s respective repository. If you have multiple
projects, you might want to arrange them in groups inside the
repository, perhaps putting projects with similar goals or shared code
in the same subdirectory, or maybe just grouping them alphabetically.
Such an arrangement might look like this:

\begin{verbatim}
/
   utils/
      calc/
         trunk/
         tags/
         branches/
      calendar/
         trunk/
         tags/
         branches/
      …
   office/
      spreadsheet/
         trunk/
         tags/
         branches/
      …
\end{verbatim}

Lay out your repository in whatever way you see fit. Subversion does not
expect or enforce a particular layout---in its eyes, a directory is a
directory is a directory. Ultimately, you should choose the repository
arrangement that meets the needs of the people who work on the projects
that live there.

In the name of full disclosure, though, we\textquotesingle ll mention
another very common layout. In this layout, the \texttt{trunk},
\texttt{tags}, and \texttt{branches} directories live in the root
directory of your repository, and your projects are in subdirectories
beneath those, like so:

\begin{verbatim}
/
   trunk/
      calc/
      calendar/
      spreadsheet/
      …
   tags/
      calc/
      calendar/
      spreadsheet/
      …
   branches/
      calc/
      calendar/
      spreadsheet/
      …
\end{verbatim}

There\textquotesingle s nothing particularly incorrect about such a
layout, but it may or may not seem as intuitive for your users.
Especially in large, multiproject situations with many users, those
users may tend to be familiar with only one or two of the projects in
the repository. But the projects-as-branch-siblings approach tends to
deemphasize project individuality and focus on the entire set of
projects as a single entity. That\textquotesingle s a social issue,
though. We like our originally suggested arrangement for purely
practical reasons---it\textquotesingle s easier to ask about (or modify,
or migrate elsewhere) the entire history of a single project when
there\textquotesingle s a single repository path that holds the entire
history---past, present, tagged, and branched---for that project and
that project alone.

\subsection{Deciding Where and How to Host Your
Repository}\label{svn.reposadmin.basics.hosting}

Before creating your Subversion repository, an obvious question
you\textquotesingle ll need to answer is where the thing is going to
live. This is strongly connected to myriad other questions involving how
the repository will be accessed (via a Subversion server or directly),
by whom (users behind your corporate firewall or the whole world out on
the open Internet), what other services you\textquotesingle ll be
providing around Subversion (repository browsing interfaces, email-based
commit notification, etc.), your data backup strategy, and so on.

We cover server choice and configuration in
\hyperref[svn.serverconfig]{Server Configuration}, but the point
we\textquotesingle d like to briefly make here is simply that the
answers to some of these other questions might have implications that
force your hand when deciding where your repository will live. For
example, certain deployment scenarios might require accessing the
repository via a remote filesystem from multiple computers, or using
multiple repositories with syncronized contents distributed
geographically to permit more performant access to that data by users
around the globe. Addressing each possible way to deploy Subversion is
both impossible and outside the scope of this book. We simply encourage
you to evaluate your options using these pages and other sources as your
reference material and to plan ahead.

\subsection{Controlling Access to Your
Repository}\label{svn.reposadmin.basics.accesscontrol}

Access control in Subversion is almost entirely managed by the
Subversion server processes. We discuss the available Subversion servers
in \hyperref[svn.serverconfig]{Server Configuration}, and explain
path-based access control specifically in
\hyperref[svn.serverconfig.pathbasedauthz]{Path-Based Authorization}. In
addition to those user-level access control questions,
you\textquotesingle ll also want to ensure that your repository is
accessible by the programs on your hosting machine which need to access
it. Consider the OS-level user and group ownership that makes sense for
your repository. Once again, the information found in
\hyperref[svn.serverconfig]{Server Configuration} should be able to help
you make these decisions.

\section{Creating and Configuring Your
Repository}\label{svn.reposadmin.create}

Earlier in this chapter (in
\hyperref[svn.reposadmin.planning]{Strategies for Repository
Deployment}), we looked at some of the important decisions that should
be made before creating and configuring your Subversion repository. Now,
we finally get to get our hands dirty! In this section,
we\textquotesingle ll see how to actually create a Subversion repository
and configure it to perform custom actions when special repository
events occur.

\subsection{Creating the
Repository}\label{svn.reposadmin.basics.creating}

Subversion repository creation is an incredibly simple task. The
\texttt{svnadmin} utility that comes with Subversion provides a
subcommand (\texttt{svnadmin\ create}) for doing just that.

\begin{verbatim}
$ # Create a repository
$ svnadmin create /var/svn/repos
$
\end{verbatim}

Assuming that the parent directory \texttt{/var/svn} exists and that you
have sufficient permissions to modify that directory, the previous
command creates a new repository in the directory
\texttt{/var/svn/repos}, and with the default filesystem data store
(FSFS). You can explicitly select the filesystem type using the
\texttt{-\/-fs-type} argument; new repositories should always use
\texttt{fsfs}.

\begin{verbatim}
$ # Create an FSFS-backed repository
$ svnadmin create --fs-type fsfs /var/svn/repos
$
\end{verbatim}

After running this simple command, you have a Subversion repository.
Depending on how users will access this new repository, you might need
to fiddle with its filesystem permissions. But since basic system
administration is rather outside the scope of this text,
we\textquotesingle ll leave further exploration of that topic as an
exercise to the reader.

The path argument to \texttt{svnadmin} is just a regular filesystem path
and not a URL like the \texttt{svn} client program uses when referring
to repositories. Both \texttt{svnadmin} and \texttt{svnlook} are
considered server-side utilities---they are used on the machine where
the repository resides to examine or modify aspects of the repository,
and are in fact unable to perform tasks across a network. A common
mistake made by Subversion newcomers is trying to pass URLs (even
``local'' \texttt{file://} ones) to these two programs.

Present in the \texttt{db/} subdirectory of your repository is the
implementation of the versioned filesystem. Your new
repository\textquotesingle s versioned filesystem begins life at
revision 0, which is defined to consist of nothing but the top-level
root (\texttt{/}) directory. Initially, revision 0 also has a single
revision property, \texttt{svn:date}, set to the time at which the
repository was created.

Now that you have a repository, it\textquotesingle s time to customize
it.

While some parts of a Subversion repository---such as the configuration
files and hook scripts---are meant to be examined and modified manually,
you shouldn\textquotesingle t (and shouldn\textquotesingle t need to)
tamper with the other parts of the repository ``by hand.'' The
\texttt{svnadmin} tool should be sufficient for any changes necessary to
your repository, or you can look to third-party tools for tweaking
relevant subsections of the repository. Do \emph{not} attempt manual
manipulation of your version control history by poking and prodding
around in your repository\textquotesingle s data store files!

\subsection{Implementing Repository Hooks}\label{svn.reposadmin.hooks}

A hook is a program triggered by some repository event, such as the
creation of a new revision or the modification of an unversioned
property. Some hooks (the so-called ``pre hooks'') run in advance of a
repository operation and provide a means by which to both report what is
about to happen and prevent it from happening at all. Other hooks (the
``post hooks'') run after the completion of a repository event and are
useful for performing tasks that examine---but don\textquotesingle t
modify---the repository. Each hook is handed enough information to tell
what that event is (or was), the specific repository changes proposed
(or completed), and the username of the person who triggered the event.

The \texttt{hooks} subdirectory is, by default, filled with templates
for various repository hooks:

\begin{verbatim}
$ ls repos/hooks/
post-commit.tmpl          post-unlock.tmpl  pre-revprop-change.tmpl
post-lock.tmpl            pre-commit.tmpl   pre-unlock.tmpl
post-revprop-change.tmpl  pre-lock.tmpl     start-commit.tmpl
$
\end{verbatim}

There is one template for each hook that the Subversion repository
supports; by examining the contents of those template scripts, you can
see what triggers each script to run and what data is passed to that
script. Also present in many of these templates are examples of how one
might use that script, in conjunction with other Subversion-supplied
programs, to perform common useful tasks. To actually install a working
hook, you need only place some executable program or script into the
\texttt{repos/hooks} directory, which can be executed as the name (such
as \texttt{start-commit} or \texttt{post-commit}) of the hook.

On Unix platforms, this means supplying a script or program (which could
be a shell script, a Python program, a compiled C binary, or any number
of other things) named exactly like the name of the hook. Of course, the
template files are present for more than just informational
purposes---the easiest way to install a hook on Unix platforms is to
simply copy the appropriate template file to a new file that lacks the
\texttt{.tmpl} extension, customize the hook\textquotesingle s contents,
and ensure that the script is executable. Windows, however, uses file
extensions to determine whether a program is executable, so you would
need to supply a program whose basename is the name of the hook and
whose extension is one of the special extensions recognized by Windows
for executable programs, such as \texttt{.exe} for programs and
\texttt{.bat} for batch files.

Subversion executes hooks as the same user who owns the process that is
accessing the Subversion repository. In most cases, the repository is
being accessed via a Subversion server, so this user is the same user as
whom the server runs on the system. The hooks themselves will need to be
configured with OS-level permissions that allow that user to execute
them. Also, this means that any programs or files (including the
Subversion repository) accessed directly or indirectly by the hook will
be accessed as the same user. In other words, be alert to potential
permission-related problems that could prevent the hook from performing
the tasks it is designed to perform.

There are several hooks implemented by the Subversion repository, and
you can get details about each of them in
\hyperref[svn.ref.reposhooks]{Subversion Repository Hook Reference}. As
a repository administrator, you\textquotesingle ll need to decide which
hooks you wish to implement (by way of providing an appropriately named
and permissioned hook program), and how. When you make this decision,
keep in mind the big picture of how your repository is deployed. For
example, if you are using server configuration to determine which users
are permitted to commit changes to your repository, you
don\textquotesingle t need to do this sort of access control via the
hook system.

\subsubsection{Hook script environment
configuration}\label{svn.reposadmin.hooks.configuration}

By default, Subversion executes hook scripts with an empty
environment---that is, no environment variables are set at all, not even
\texttt{\$PATH} (or \texttt{\%PATH\%}, under Windows). Because of this,
many administrators are baffled when their hook program runs fine by
hand, but doesn\textquotesingle t work when invoked by Subversion.
Administrators have historically worked around this problem by manually
setting all the environment variables their hook scripts need in the
scripts themselves.

Subversion 1.8 introduces a new way to manage the environment of
Subversion-executed hook scripts---the hook script environment
configuration file. If a Subversion server finds a file named
\texttt{hooks-env} in the repository\textquotesingle s \texttt{conf/}
subdirectory, it parses that file as an INI-formatted configuration file
and applies the option names and variables found therein to the hook
script\textquotesingle s execution environment as environment variables.

The syntax of the \texttt{hooks-env} file is pretty straightforward:
each section name is the name of a hook script (such as
\texttt{{[}pre-commit{]}} or \texttt{{[}post-revprop-change{]}}), and
the configuration items inside that section are treated as mappings of
environment variable names to desired values. Additionally, there is a
special \texttt{{[}default{]}} section, which can be used to configure
environment variable mappings that should be applied to \emph{all} hook
scripts (unless explicitly overridden by per-hook-script settings). See
\hyperref[svn.reposadmin.hooks.configuration.ex-1]{hooks-env (custom
hook script environment configuration)} for a sample \texttt{hooks-env}
configuration file.

\protect\phantomsection\label{svn.reposadmin.hooks.configuration.ex-1}
hooks-env (custom hook script environment configuration)

\begin{verbatim}
# All scripts should use a UTF-8 locale and have our hook script
# utilities directory on the search path.

[default]
LANG = en_US.UTF-8
PATH = /usr/local/svn/tools:/usr/bin


# The post-commit and post-revprop-change scripts want to run
# programs from our custom synctools replication software suite, too.

[post-commit]
PATH = /usr/local/synctools-1.1/bin:%(PATH)s

[post-revprop-change]
PATH = /usr/local/synctools-1.1/bin:%(PATH)s
\end{verbatim}

\hyperref[svn.reposadmin.hooks.configuration.ex-1]{hooks-env (custom
hook script environment configuration)} also demonstrates the nifty
string substitution syntax found in Subversion\textquotesingle s
configuration file parser. In this example, the value of the
\texttt{PATH} option---pulled from the \texttt{{[}default{]}} section of
the file---is substituted in place of the \texttt{\%(PATH)s} placeholder
text in the per-hook sections. For more about this special syntax, see
the \texttt{README.txt} file which lives in the Subversion runtime
configuration directory. (And for more information about that directory,
see \hyperref[svn.advanced.confarea]{Runtime Configuration Area}.)

Of course, having exact duplicates of your custom hook script
environment configuration files in every single
repository\textquotesingle s \texttt{conf/} directory could get
cumbersome, especially when you need to make changes to them all. So
Subversion\textquotesingle s servers allow you to specify an alternate
(possibly shared) location for this configuration information.

\subsubsection{Common uses for hook
scripts}\label{svn.reposadmin.hooks.uses}

Repository hook scripts can offer a wide range of utility, but most tend
to fall into a few basic categories: notification, validation, and
replication.

Notification scripts are those which tell someone that something
happened. The most common of these found in a Subversion service
offering involve programs which send commit and revision property change
notification emails to project members, driven by the post-commit and
post-revprop-change hooks, respectively. There are numerous other
notification approaches, from issue tracker integration scripts to
scripts which operate as IRC bots to announce that
something\textquotesingle s changed in the repository.

On the validation side of things, the start-commit and pre-commit hooks
are widely used to allow or disallow commits based on various criteria:
the author of the commit, the formatting and/or content of the log
message which describes the commit, and even the low-level details of
the changes made to files and directories in the commit. Likewise, the
pre-revprop-change hook acts as the gateway to revision property
changes, which is an especially valuable role considering the fact that
revision properties are not themselves versioned, and can therefore only
be modified destructively.

One special class of change validation that has seen widespread use
since Subversion 1.5 was released is validation of the committing client
software itself. When Subversion\textquotesingle s merge tracking
feature (described extensively in \hyperref[svn.branchmerge]{Branching
and Merging}) was introduced in that release, Subversion administrators
needed a way to ensure that once users of their repositories started
using the new feature that \emph{all} their merges were tracked. To
reduce the chance of someone committing an untracked merge to the
repository, they used start-commit hooks to examine the feature
capabilities string advertised by Subversion clients. If the committing
client didn\textquotesingle t advertise support for merge tracking, the
commit was denied with instructions to the user to immediately update
their Subversion client!
\hyperref[svn.reposadmin.hooks.uses.ex-1]{start-commit hook to require
merge tracking support} provides an example of a start-commit script
which does precisely this.

\protect\phantomsection\label{svn.reposadmin.hooks.uses.ex-1}
start-commit hook to require merge tracking support

\begin{verbatim}
#!/usr/bin/env python
import sys

# sys.argv[3] is a colon-delimited capabilities list
if 'mergeinfo' not in sys.argv[3].split(':'):
  sys.stderr.write("""\
ERROR: Commits to this repository must be made using Subversion
clients which support the merge tracking feature.  Please upgrade
your client to at least Subversion 1.5.0.
""")
  sys.exit(1)
\end{verbatim}

Beginning in Subversion 1.8, clients committing against a Subversion 1.8
server will still provide the feature capabilities string, but will also
provide additional information about themselves by way of ephemeral
transaction properties. Ephemeral transaction properties are essentially
revision properties which are set on the commit transaction by the
client at the earliest opportunity while committing, but which are
automatically removed by the server immediately prior to the transaction
becoming a finalized revision. You can inspect these properties using
the same tools with which you\textquotesingle d inspect other
unversioned properties set on commit transactions during the timeframe
between which the start-commit and pre-commit repository hook scripts
would operate.

The following are the ephemeral transaction properties which Subversion
currently provides and implements:

\begin{description}
\item[\texttt{svn:txn-client-compat-version}]
Carries the Subversion library version string with which the committing
client claims compatibility. This is useful for deciding whether the
client supports the minimal feature set required for proper handling of
the repository data.
\item[\texttt{svn:txn-user-agent}]
Carries the ``user agent'' string which describes the committing client
program. Subversion\textquotesingle s libraries define the initial
portion of this string, but third-party consumers of the API (GUI
clients, etc.) can append custom information to it.
\end{description}

While most clients will transmit ephemeral transaction properties early
enough in the commit process that they may be inspected by the
start-commit hook script, some configurations of Subversion will cause
those properties to not be set on the transaction until later in the
commit process. Administrators should consider performing any validation
based on ephemeral transaction properties in both the start-commit and
pre-commit hooks---the former to rule out invalid clients before those
clients transmit the commit payload; the latter ``just in case'' the
validation checks couldn\textquotesingle t be performed by the
start-commit hook.

As noted before, ephemeral transaction properties are removed from the
transaction just before it is promoted to a new revision. Some
administrators may wish to preserve the information in those properties
indefinitely. We suggest that you do so by using the pre-commit hook
script to copy the values of those properties to new property names. In
fact, the Subversion source code distribution provides a
\texttt{persist-ephemeral-txnprops.py} script (in the
\texttt{tools/hook-scripts/} subdirectory) for doing precisely that.

The third common type of hook script usage is for the purpose of
replication. Whether you are driving a simple backup process or a more
involved remote repository mirroring scenario, hook scripts can be
critical. See \hyperref[svn.reposadmin.maint.backup]{Repository Backup}
and \hyperref[svn.reposadmin.maint.replication]{Repository Replication}
for more information about these aspects of repository maintenance.

\subsubsection{Finding hook scripts or rolling your
own}\label{svn.reposadmin.hooks.summary}

As you might imagine, there is no shortage of Subversion hook programs
and scripts that are freely available either from the Subversion
community itself or elsewhere. In fact, the Subversion distribution
provides several commonly used hook scripts in its
\texttt{tools/hook-scripts/} subdirectory. However, if you are unable to
find one that meets your specific needs, you might consider writing your
own. See \hyperref[svn.developer]{Embedding Subversion} for information
about developing software using Subversion\textquotesingle s public
APIs.

Hook scripts can do almost anything, but hook script authors should show
restraint. It might be tempting to, say, use hook scripts to
automatically correct errors, shortcomings, or policy violations present
in the files being committed. Unfortunately, doing so can cause
problems. Subversion keeps client-side caches of certain bits of
repository data, and if you change a commit transaction in this way,
those caches become indetectably stale, leading to surprising and
unexpected behavior. While it is generally okay to add new commit
transaction properties via a hook script, essentially everything else
about a commit transaction should be considered read-only. Instead of
modifying a transaction to polish its payload, simply \emph{validate}
the transaction in the pre-commit hook and reject the commit if it does
not meet the desired requirements. As a bonus, your users will learn the
value of careful, compliance-minded work habits.

\subsection{FSFS Configuration}\label{svn.reposadmin.create.fsfs}

As of Subversion 1.6, FSFS filesystems have several configurable
parameters which an administrator can use to fine-tune the performance
or disk usage of their repositories. You can find these options---and
the documentation for them---in the \texttt{db/fsfs.conf} file in the
repository.

The on-disk FSFS format has evolved over time, and a newly created
repository uses the most capable format its version of Subversion
supports. Format 7, introduced in Subversion 1.9, switched FSFS to
\emph{logical addressing} and added full checksum coverage,
substantially reducing the amount of disk I/O needed to read revision
data. Format 8, introduced in Subversion 1.10, added support for LZ4
compression. (If a new repository must remain compatible with an older
server, create it with the \texttt{-\/-compatible-version} option to
\texttt{svnadmin\ create} so that it uses an older format.)

The settings most worth knowing about are the following:

\begin{description}
\item[\texttt{compression}]
Selects the algorithm used to compress file contents:
\texttt{lz4}, \texttt{zlib}, \texttt{zlib-1} through \texttt{zlib-9}, or
\texttt{none}. LZ4---available only on format 8 and later
repositories---is very fast and is the default for such repositories;
zlib compresses more tightly but more slowly, and a bare \texttt{zlib} is
equivalent to \texttt{zlib-5}. The older \texttt{compression-level}
option, which configured the zlib level, is still honored for backward
compatibility.
\item[\texttt{enable-rep-sharing}]
When enabled (the default), identical file and property contents are
stored only once across the entire repository, even when they appear in
otherwise unrelated files. This is the data tracked by the representation
cache (see \hyperref[svn.ref.svnadmin.c.build-repcache]{svnadmin
build-repcache}).
\item[\texttt{enable-dir-deltification}, \texttt{enable-props-deltification}]
Control whether directory listings and property lists, respectively, are
stored as deltas against earlier revisions to save space. Both are
enabled by default.
\item[\texttt{max-deltification-walk}, \texttt{max-linear-deltification}]
Tune how aggressively Subversion deltifies stored representations,
trading disk space against the work required to reconstruct full text.
\end{description}

Changing these options affects only revisions written after the change;
existing data is left as it is until it is rewritten---for example, by a
dump and load.

\section{Repository Maintenance}\label{svn.reposadmin.maint}

Maintaining a Subversion repository can be daunting, mostly due to the
complexities inherent in systems that have a database backend. Doing the
task well is all about knowing the tools---what they are, when to use
them, and how. This section will introduce you to the repository
administration tools provided by Subversion and discuss how to wield
them to accomplish tasks such as repository data migration, upgrades,
backups, and cleanups.

\subsection{An Administrator\textquotesingle s
Toolkit}\label{svn.reposadmin.maint.tk}

Subversion provides a handful of utilities useful for creating,
inspecting, modifying, and repairing your repository.
Let\textquotesingle s look more closely at each of those tools.

\subsubsection{svnadmin}\label{svn.reposadmin.maint.tk.svnadmin}

The \texttt{svnadmin} program is the repository
administrator\textquotesingle s best friend. Besides providing the
ability to create Subversion repositories, this program allows you to
perform several maintenance operations on those repositories. The syntax
of \texttt{svnadmin} is similar to that of other Subversion command-line
programs:

\begin{verbatim}
$ svnadmin help
general usage: svnadmin SUBCOMMAND REPOS_PATH  [ARGS & OPTIONS ...]
Type 'svnadmin help <subcommand>' for help on a specific subcommand.
Type 'svnadmin --version' to see the program version and FS modules.

Available subcommands:
   build-repcache
   crashtest
   create
…
\end{verbatim}

Previously in this chapter (in
\hyperref[svn.reposadmin.basics.creating]{Creating the Repository}), we
were introduced to the \texttt{svnadmin\ create} subcommand. Most of the
other \texttt{svnadmin} subcommands we will cover later in this chapter.
And you can consult \hyperref[svn.ref.svnadmin]{svnadmin
Reference---Subversion Repository Administration} for a full rundown of
subcommands and what each of them offers.

\subsubsection{svnlook}\label{svn.reposadmin.maint.tk.svnlook}

\texttt{svnlook} is a tool provided by Subversion for examining the
various revisions and transactions (which are revisions in the making)
in a repository. No part of this program attempts to change the
repository. \texttt{svnlook} is typically used by the repository hooks
for reporting the changes that are about to be committed (in the case of
the pre-commit hook) or that were just committed (in the case of the
post-commit hook) to the repository. A repository administrator may use
this tool for diagnostic purposes.

\texttt{svnlook} has a straightforward syntax:

\begin{verbatim}
$ svnlook help
general usage: svnlook SUBCOMMAND REPOS_PATH [ARGS & OPTIONS ...]
Note: any subcommand which takes the '--revision' and '--transaction'
      options will, if invoked without one of those options, act on
      the repository's youngest revision.
Type 'svnlook help <subcommand>' for help on a specific subcommand.
Type 'svnlook --version' to see the program version and FS modules.
…
\end{verbatim}

Most of \texttt{svnlook}\textquotesingle s subcommands can operate on
either a revision or a transaction tree, printing information about the
tree itself, or how it differs from the previous revision of the
repository. You use the \texttt{-\/-revision} (\texttt{-r}) and
\texttt{-\/-transaction} (\texttt{-t}) options to specify which revision
or transaction, respectively, to examine. In the absence of both the
\texttt{-\/-revision} (\texttt{-r}) and \texttt{-\/-transaction}
(\texttt{-t}) options, \texttt{svnlook} will examine the youngest (or
\texttt{HEAD}) revision in the repository. So the following two commands
do exactly the same thing when 19 is the youngest revision in the
repository located at \texttt{/var/svn/repos}:

\begin{verbatim}
$ svnlook info /var/svn/repos
$ svnlook info /var/svn/repos -r 19
\end{verbatim}

One exception to these rules about subcommands is the
\texttt{svnlook\ youngest} subcommand, which takes no options and simply
prints out the repository\textquotesingle s youngest revision number:

\begin{verbatim}
$ svnlook youngest /var/svn/repos
19
$
\end{verbatim}

Keep in mind that the only transactions you can browse are uncommitted
ones. Most repositories will have no such transactions because
transactions are usually either committed (in which case, you should
access them as revision with the \texttt{-\/-revision} (\texttt{-r})
option) or aborted and removed.

Output from \texttt{svnlook} is designed to be both human- and
machine-parsable. Take, as an example, the output of the
\texttt{svnlook\ info} subcommand:

\begin{verbatim}
$ svnlook info /var/svn/repos
sally
2002-11-04 09:29:13 -0600 (Mon, 04 Nov 2002)
27
Added the usual
Greek tree.
$
\end{verbatim}

The output of \texttt{svnlook\ info} consists of the following, in the
order given:

\begin{enumerate}
\def\labelenumi{\arabic{enumi}.}
\item
  The author, followed by a newline
\item
  The date, followed by a newline
\item
  The number of characters in the log message, followed by a newline
\item
  The log message itself, followed by a newline
\end{enumerate}

This output is human-readable, meaning items such as the datestamp are
displayed using a textual representation instead of something more
obscure (such as the number of nanoseconds since the Tastee Freez guy
drove by). But the output is also machine-parsable---because the log
message can contain multiple lines and be unbounded in length,
\texttt{svnlook} provides the length of that message before the message
itself. This allows scripts and other wrappers around this command to
make intelligent decisions about the log message, such as how much
memory to allocate for the message, or at least how many bytes to skip
in the event that this output is not the last bit of data in the stream.

\texttt{svnlook} can perform a variety of other queries: displaying
subsets of bits of information we\textquotesingle ve mentioned
previously, recursively listing versioned directory trees, reporting
which paths were modified in a given revision or transaction, showing
textual and property differences made to files and directories, and so
on. See \hyperref[svn.ref.svnlook]{svnlook Reference---Subversion
Repository Examination} for a full reference of
\texttt{svnlook}\textquotesingle s features.

\subsubsection{svndumpfilter}\label{svn.reposadmin.maint.tk.svndumpfilter}

While it won\textquotesingle t be the most commonly used tool at the
administrator\textquotesingle s disposal, \texttt{svndumpfilter}
provides a very particular brand of useful functionality---the ability
to quickly and easily modify streams of Subversion repository history
data by acting as a path-based filter.

The syntax of \texttt{svndumpfilter} is as follows:

\begin{verbatim}
$ svndumpfilter help
general usage: svndumpfilter SUBCOMMAND [ARGS & OPTIONS ...]
Type 'svndumpfilter help <subcommand>' for help on a specific subcommand.
Type 'svndumpfilter --version' to see the program version.
  
Available subcommands:
   exclude
   include
   help (?, h)
\end{verbatim}

There are only two interesting subcommands:
\texttt{svndumpfilter\ exclude} and \texttt{svndumpfilter\ include}.
They allow you to make the choice between implicit or explicit inclusion
of paths in the stream. You can learn more about these subcommands and
\texttt{svndumpfilter}\textquotesingle s unique purpose later in this
chapter, in \hyperref[svn.reposadmin.maint.filtering]{Filtering
Repository History}.

\subsubsection{svnrdump}\label{svn.reposadmin.maint.tk.svnrdump}

The \texttt{svnrdump} program is, to put it simply, essentially just
network-aware flavors of the \texttt{svnadmin\ dump} and
\texttt{svnadmin\ load} subcommands, rolled up into a separate program.

\begin{verbatim}
$ svnrdump help
general usage: svnrdump SUBCOMMAND URL [-r LOWER[:UPPER]]
Type 'svnrdump help <subcommand>' for help on a specific subcommand.
Type 'svnrdump --version' to see the program version and RA modules.

Available subcommands:
   dump
   load
   help (?, h)

$
\end{verbatim}

We discuss the use of \texttt{svnrdump} and the aforementioned
\texttt{svnadmin} commands later in this chapter (see
\hyperref[svn.reposadmin.maint.migrate]{Migrating Repository Data
Elsewhere}).

\subsubsection{svnsync}\label{svn.reposadmin.maint.tk.svnsync}

The \texttt{svnsync} program provides all the functionality required for
maintaining a read-only mirror of a Subversion repository. The program
really has one job---to transfer one repository\textquotesingle s
versioned history into another repository. And while there are few ways
to do that, its primary strength is that it can operate remotely---the
``source'' and ``sink''\footnote{Or is that, the ``sync''?} repositories
may be on different computers from each other and from \texttt{svnsync}
itself.

As you might expect, \texttt{svnsync} has a syntax that looks very much
like every other program we\textquotesingle ve mentioned in this
chapter:

\begin{verbatim}
$ svnsync help
general usage: svnsync SUBCOMMAND DEST_URL  [ARGS & OPTIONS ...]
Type 'svnsync help <subcommand>' for help on a specific subcommand.
Type 'svnsync --version' to see the program version and RA modules.

Available subcommands:
   initialize (init)
   synchronize (sync)
   copy-revprops
   info
   help (?, h)
$
\end{verbatim}

We talk more about replicating repositories with \texttt{svnsync} later
in this chapter (see
\hyperref[svn.reposadmin.maint.replication]{Repository Replication}).

\subsubsection{fsfs-reshard.py}\label{svn.reposadmin.maint.tk.fsfsreshard}

While not an official member of the Subversion toolchain, the
\texttt{fsfs-reshard.py} script (found in the \texttt{tools/server-side}
directory of the Subversion source distribution) is a useful performance
tuning tool for administrators of FSFS-backed Subversion repositories.
FSFS repositories use individual files to house information about each
revision. Sometimes these files all live in a single directory;
sometimes they are sharded across many directories.

The earliest FSFS release versions would house all the revision files
within a single directory that grew---one file per revision---throughout
the lifetime of your repository. This created problems on systems which
have hard limits on the number of files permitted in a given directory,
and was a performance burden even on systems where such limits
didn\textquotesingle t exist or were set sufficiently high.

Beginning in version 1.5, Subversion creates FSFS-backed repositories
using a slightly modified layout in which the contents of the revision
files directory (and other always-growing directories) are sharded, or
scattered across several subdirectories. This can greatly reduce the
time it takes the system to locate any one of these files, and therefore
increases the overall performance of Subversion when reading from the
repository.

The number of files permitted to live in a given subdirectory is a
configurable thing (though the defaults are reasonable ones for most
known platforms), but changing that configuration after the repository
has been in use for some time could cause Subversion to be unable to
locate the files it is looking for. That\textquotesingle s where
\texttt{fsfs-reshard.py} comes in.

\texttt{fsfs-reshard.py} reshuffles the repository\textquotesingle s
file structure into a new arrangement that reflects the requested number
of sharding subdirectories and updates the repository configuration to
preserve this change. When used in conjunction with the
\texttt{svnadmin\ upgrade} command, this is especially useful for
upgrading a pre-1.5 Subversion (unsharded) repository to the latest
filesystem format and sharding its data files (which Subversion will not
automatically do for you). This script can also be used for fine-tuning
an already sharded repository.

\subsection{Commit Log Message
Correction}\label{svn.reposadmin.maint.setlog}

Sometimes a user will have an error in her log message (a misspelling or
some misinformation, perhaps). If the repository is configured (using
the pre-revprop-change hook; see
\hyperref[svn.reposadmin.hooks]{Implementing Repository Hooks}) to
accept changes to this log message after the commit is finished, the
user can ``fix'' her log message remotely using \texttt{svn\ propset}
(see \hyperref[svn.ref.svn.c.propset]{svn propset (pset, ps)} in
\hyperref[svn.ref.svn]{svn Reference---Subversion Command-Line Client}).
However, because of the potential to lose information forever,
Subversion repositories are not, by default, configured to allow changes
to unversioned properties---except by an administrator.

If a log message needs to be changed by an administrator, this can be
done using \texttt{svnadmin\ setlog}. This command changes the log
message (the \texttt{svn:log} property) on a given revision of a
repository, reading the new value from a provided file.

\begin{verbatim}
$ echo "Here is the new, correct log message" > newlog.txt
$ svnadmin setlog myrepos newlog.txt -r 388
\end{verbatim}

The \texttt{svnadmin\ setlog} command, by default, is still bound by the
same protections against modifying unversioned properties as a remote
client is---the pre-revprop-change and post-revprop-change hooks are
still triggered, and therefore must be set up to accept changes of this
nature. But an administrator can get around these protections by passing
the \texttt{-\/-bypass-hooks} option to the \texttt{svnadmin\ setlog}
command.

Remember, though, that by bypassing the hooks, you are likely avoiding
such things as email notifications of property changes, backup systems
that track unversioned property changes, and so on. In other words, be
very careful about what you are changing, and how you change it.

\subsection{Managing Disk Space}\label{svn.reposadmin.maint.diskspace}

While the cost of storage has dropped incredibly in the past few years,
disk usage is still a valid concern for administrators seeking to
version large amounts of data. Every bit of version history information
stored in the live repository needs to be backed up elsewhere, perhaps
multiple times as part of rotating backup schedules. It is useful to
know what pieces of Subversion\textquotesingle s repository data need to
remain on the live site, which need to be backed up, and which can be
safely removed.

\subsubsection{How Subversion saves disk
space}\label{svn.reposadmin.maint.diskspace.deltas}

To keep the repository small, Subversion uses deltification (or
delta-based storage) within the repository itself. Deltification
involves encoding the representation of a chunk of data as a collection
of differences against some other chunk of data. If the two pieces of
data are very similar, this deltification results in storage savings for
the deltified chunk---rather than taking up space equal to the size of
the original data, it takes up only enough space to say, ``I look just
like this other piece of data over here, except for the following couple
of changes.'' The result is that most of the repository data that tends
to be bulky---namely, the contents of versioned files---is stored at a
much smaller size than the original full-text representation of that
data.

While deltified storage has been a part of Subversion\textquotesingle s
design since the very beginning, there have been additional improvements
made over the years. Subversion repositories created with Subversion 1.4
or later benefit from compression of the full-text representations of
file contents. Repositories created with Subversion 1.6 or later further
enjoy the disk space savings afforded by representation sharing, a
feature which allows multiple files or file revisions with identical
file content to refer to a single shared instance of that data rather
than each having their own distinct copy thereof.

\subsubsection{Removing dead
transactions}\label{svn.reposadmin.maint.diskspace.deadtxns}

Though they are uncommon, there are circumstances in which a Subversion
commit process might fail, leaving behind in the repository the remnants
of the revision-to-be that wasn\textquotesingle t---an uncommitted
transaction and all the file and directory changes associated with it.
This could happen for several reasons: perhaps the client operation was
inelegantly terminated by the user, or a network failure occurred in the
middle of an operation. Regardless of the reason, dead transactions can
happen. They don\textquotesingle t do any real harm, other than
consuming disk space. A fastidious administrator may nonetheless wish to
remove them.

You can use the \texttt{svnadmin\ lstxns} command to list the names of
the currently outstanding transactions:

\begin{verbatim}
$ svnadmin lstxns myrepos
19
3a1
a45
$
\end{verbatim}

Each item in the resultant output can then be used with \texttt{svnlook}
(and its \texttt{-\/-transaction} (\texttt{-t}) option) to determine who
created the transaction, when it was created, what types of changes were
made in the transaction---information that is helpful in determining
whether the transaction is a safe candidate for removal! If you do
indeed want to remove a transaction, its name can be passed to
\texttt{svnadmin\ rmtxns}, which will perform the cleanup of the
transaction. In fact, \texttt{svnadmin\ rmtxns} can take its input
directly from the output of \texttt{svnadmin\ lstxns}!

\begin{verbatim}
$ svnadmin rmtxns myrepos `svnadmin lstxns myrepos`
$
\end{verbatim}

If you use these two subcommands like this, you should consider making
your repository temporarily inaccessible to clients. That way, no one
can begin a legitimate transaction before you start your cleanup.
\hyperref[svn.reposadmin.maint.diskspace.deadtxns.ex-1]{txn-info.sh
(reporting outstanding transactions)} contains a bit of shell-scripting
that can quickly generate information about each outstanding transaction
in your repository.

\protect\phantomsection\label{svn.reposadmin.maint.diskspace.deadtxns.ex-1}
txn-info.sh (reporting outstanding transactions)

\begin{verbatim}
#!/bin/sh

### Generate informational output for all outstanding transactions in
### a Subversion repository.

REPOS="${1}"
if [ "x$REPOS" = x ] ; then
  echo "usage: $0 REPOS_PATH"
  exit
fi

for TXN in `svnadmin lstxns ${REPOS}`; do 
  echo "---[ Transaction ${TXN} ]-------------------------------------------"
  svnlook info "${REPOS}" -t "${TXN}"
done
\end{verbatim}

The output of the script is basically a concatenation of several chunks
of \texttt{svnlook\ info} output (see
\hyperref[svn.reposadmin.maint.tk.svnlook]{svnlook}) and will look
something like this:

\begin{verbatim}
$ txn-info.sh myrepos
---[ Transaction 19 ]-------------------------------------------
sally
2001-09-04 11:57:19 -0500 (Tue, 04 Sep 2001)
0
---[ Transaction 3a1 ]-------------------------------------------
harry
2001-09-10 16:50:30 -0500 (Mon, 10 Sep 2001)
39
Trying to commit over a faulty network.
---[ Transaction a45 ]-------------------------------------------
sally
2001-09-12 11:09:28 -0500 (Wed, 12 Sep 2001)
0
$
\end{verbatim}

A long-abandoned transaction usually represents some sort of failed or
interrupted commit. A transaction\textquotesingle s datestamp can
provide interesting information---for example, how likely is it that an
operation begun nine months ago is still active?

In short, transaction cleanup decisions need not be made unwisely.
Various sources of information---including Apache\textquotesingle s
error and access logs, Subversion\textquotesingle s operational logs,
Subversion revision history, and so on---can be employed in the
decision-making process. And of course, an administrator can often
simply communicate with a seemingly dead transaction\textquotesingle s
owner (via email, e.g.) to verify that the transaction is, in fact, in a
zombie state.

\subsubsection{Packing FSFS
filesystems}\label{svn.reposadmin.maint.diskspace.fsfspacking}

FSFS repositories contain files that describe the changes made in a
single revision, and files that contain the revision properties
associated with a single revision. Repositories created in versions of
Subversion prior to 1.5 keep these files in two directories---one for
each type of file. As new revisions are committed to the repository,
Subversion drops more files into these two directories---over time, the
number of these files in each directory can grow to be quite large. This
has been observed to cause performance problems on certain network-based
filesystems.

The first problem is that the operating system has to reference many
different files over a short period of time. This leads to inefficient
use of disk caches and, as a result, more time spent seeking across
large disks. Because of this, Subversion pays a performance penalty when
accessing your versioned data.

The second problem is a bit more subtle. Because of the ways that most
filesystems allocate disk space, each file claims more space on the disk
than it actually uses. The amount of extra space required to house a
single file can average anywhere from 2 to 16 kilobytes \emph{per file},
depending on the underlying filesystem in use. This translates directly
into a per-revision disk usage penalty for FSFS-backed repositories. The
effect is most pronounced in repositories which have many small
revisions, since the overhead involved in storing the revision file
quickly outgrows the size of the actual data being stored.

To solve these problems, Subversion 1.6 introduced the
\texttt{svnadmin\ pack} command. By concatenating all the files of a
completed shard into a single ``pack'' file and then removing the
original per-revision files, \texttt{svnadmin\ pack} reduces the file
count within a given shard down to just a single file. In doing so, it
aids filesystem caches and reduces (to one) the number of times a file
storage overhead penalty is paid.

Subversion can pack existing sharded repositories which have been
upgraded to the 1.6 filesystem format or later (see
\hyperref[svn.ref.svnadmin.c.upgrade]{svnadmin upgrade}) in
\hyperref[svn.ref.svnadmin]{svnadmin Reference---Subversion Repository
Administration}. To do so, just run \texttt{svnadmin\ pack} on the
repository:

\begin{verbatim}
$ svnadmin pack /var/svn/repos
Packing shard 0...done.
Packing shard 1...done.
Packing shard 2...done.
…
Packing shard 34...done.
Packing shard 35...done.
Packing shard 36...done.
$
\end{verbatim}

Because the packing process obtains the required locks before doing its
work, you can run it on live repositories, or even as part of a
post-commit hook. Repacking packed shards is legal, but will have no
effect on the disk usage of the repository.

\subsection{Migrating Repository Data
Elsewhere}\label{svn.reposadmin.maint.migrate}

A Subversion filesystem has its data spread throughout files in the
repository, in a fashion generally understood by (and of interest to)
only the Subversion developers themselves. However, circumstances may
arise that call for all, or some subset, of that data to be copied or
moved into another repository.

Subversion provides such functionality by way of repository dump
streams. A repository dump stream (often referred to as a ``dump file''
when stored as a file on disk) is a portable, flat file format that
describes the various revisions in your repository---what was changed,
by whom, when, and so on. This dump stream is the primary mechanism used
to marshal versioned history---in whole or in part, with or without
modification---between repositories. And Subversion provides the tools
necessary for creating and loading these dump streams: the
\texttt{svnadmin\ dump} and \texttt{svnadmin\ load} subcommands,
respectively, and the \texttt{svnrdump} program.

While the Subversion repository dump format contains human-readable
portions and a familiar structure (it resembles an RFC 822 format, the
same type of format used for most email), it is \emph{not} a plain-text
file format. It is a binary file format, highly sensitive to meddling.
For example, many text editors will corrupt the file by automatically
converting line endings.

There are many reasons for dumping and loading Subversion repository
data. Early in Subversion\textquotesingle s life, the most common reason
was due to the evolution of Subversion itself. As Subversion matured,
there were times when changes made to the backend database schema caused
compatibility issues with previous versions of the repository, so users
had to dump their repository data using the previous version of
Subversion and load it into a freshly created repository with the new
version of Subversion. Now, these types of schema changes
haven\textquotesingle t occurred since Subversion\textquotesingle s 1.0
release, and the Subversion developers promise not to force users to
dump and load their repositories when upgrading between minor versions
(such as from 1.3 to 1.4) of Subversion. But there are still other
reasons for dumping and loading, such as migrating repository history
into another repository, or (as we\textquotesingle ll cover later
in this chapter in \hyperref[svn.reposadmin.maint.filtering]{Filtering
Repository History}) purging versioned data from repository history.

The Subversion repository dump format describes versioned repository
changes only. It will not carry any information about uncommitted
transactions, user locks on filesystem paths, repository or server
configuration customizations (including hook scripts), and so on.

The Subversion repository dump format also enables conversion from a
different storage mechanism or version control system altogether.
Because the dump file format is, for the most part, human-readable, it
should be relatively easy to describe generic sets of changes---each of
which should be treated as a new revision---using this file format. In
fact, the \texttt{cvs2svn} utility uses the dump format to represent the
contents of a CVS repository so that those contents can be copied into a
Subversion repository.

For now, we\textquotesingle ll concern ourselves only with migration of
repository data between Subversion repositories, which
we\textquotesingle ll describe in detail in the sections which follow.

\subsubsection{Repository data migration using
svnadmin}\label{svn.reposadmin.maint.migrate.svnadmin}

Whatever your reason for migrating repository history, using the
\texttt{svnadmin\ dump} and \texttt{svnadmin\ load} subcommands is
straightforward. \texttt{svnadmin\ dump} will output a range of
repository revisions that are formatted using
Subversion\textquotesingle s custom filesystem dump format. The dump
format is printed to the standard output stream, while informative
messages are printed to the standard error stream. This allows you to
redirect the output stream to a file while watching the status output in
your terminal window. For example:

\begin{verbatim}
$ svnlook youngest myrepos
26
$ svnadmin dump myrepos > dumpfile
* Dumped revision 0.
* Dumped revision 1.
* Dumped revision 2.
…
* Dumped revision 25.
* Dumped revision 26.
\end{verbatim}

At the end of the process, you will have a single file
(\texttt{dumpfile} in the previous example) that contains all the data
stored in your repository in the requested range of revisions. Note that
\texttt{svnadmin\ dump} is reading revision trees from the repository
just like any other ``reader'' process would (e.g.,
\texttt{svn\ checkout}), so it\textquotesingle s safe to run this
command at any time.

The other subcommand in the pair, \texttt{svnadmin\ load}, parses the
standard input stream as a Subversion repository dump file and
effectively replays those dumped revisions into the target repository
for that operation. It also gives informative feedback, this time using
the standard output stream:

\begin{verbatim}
$ svnadmin load newrepos < dumpfile
<<< Started new txn, based on original revision 1
     * adding path : A ... done.
     * adding path : A/B ... done.
     …
------- Committed new rev 1 (loaded from original rev 1) >>>

<<< Started new txn, based on original revision 2
     * editing path : A/mu ... done.
     * editing path : A/D/G/rho ... done.

------- Committed new rev 2 (loaded from original rev 2) >>>

…

<<< Started new txn, based on original revision 25
     * editing path : A/D/gamma ... done.

------- Committed new rev 25 (loaded from original rev 25) >>>

<<< Started new txn, based on original revision 26
     * adding path : A/Z/zeta ... done.
     * editing path : A/mu ... done.

------- Committed new rev 26 (loaded from original rev 26) >>>
\end{verbatim}

The result of a load is new revisions added to a repository---the same
thing you get by making commits against that repository from a regular
Subversion client. Just as in a commit, you can use hook programs to
perform actions before and after each of the commits made during a load
process. By passing the \texttt{-\/-use-pre-commit-hook} and
\texttt{-\/-use-post-commit-hook} options to \texttt{svnadmin\ load},
you can instruct Subversion to execute the pre-commit and post-commit
hook programs, respectively, for each loaded revision. You might use
these, for example, to ensure that loaded revisions pass through the
same validation steps that regular commits pass through. Of course, you
should use these options with care---if your post-commit hook sends
emails to a mailing list for each new commit, you might not want to spew
hundreds or thousands of commit emails in rapid succession at that list!
You can read more about the use of hook scripts in
\hyperref[svn.reposadmin.hooks]{Implementing Repository Hooks}.

Note that because \texttt{svnadmin} uses standard input and output
streams for the repository dump and load processes, people who are
feeling especially saucy can try things such as this (perhaps even using
different versions of \texttt{svnadmin} on each side of the pipe):

\begin{verbatim}
$ svnadmin create newrepos
$ svnadmin dump oldrepos | svnadmin load newrepos
\end{verbatim}

By default, the dump file will be quite large---much larger than the
repository itself. That\textquotesingle s because by default every
version of every file is expressed as a full text in the dump file. This
is the fastest and simplest behavior, and it\textquotesingle s nice if
you\textquotesingle re piping the dump data directly into some other
process (such as a compression program, filtering program, or loading
process). But if you\textquotesingle re creating a dump file for
longer-term storage, you\textquotesingle ll likely want to save disk
space by using the \texttt{-\/-deltas} option. With this option,
successive revisions of files will be output as compressed, binary
differences---just as file revisions are stored in a repository. This
option is slower, but it results in a dump file much closer in size to
the original repository.

We mentioned previously that \texttt{svnadmin\ dump} outputs a range of
revisions. Use the \texttt{-\/-revision} (\texttt{-r}) option to specify
a single revision, or a range of revisions, to dump. If you omit this
option, all the existing repository revisions will be dumped.

\begin{verbatim}
$ svnadmin dump myrepos -r 23 > rev-23.dumpfile
$ svnadmin dump myrepos -r 100:200 > revs-100-200.dumpfile
\end{verbatim}

As Subversion dumps each new revision, it outputs only enough
information to allow a future loader to re-create that revision based on
the previous one. In other words, for any given revision in the dump
file, only the items that were changed in that revision will appear in
the dump. The only exception to this rule is the first revision that is
dumped with the current \texttt{svnadmin\ dump} command.

By default, Subversion will not express the first dumped revision as
merely differences to be applied to the previous revision. For one
thing, there is no previous revision in the dump file! And second,
Subversion cannot know the state of the repository into which the dump
data will be loaded (if it ever is). To ensure that the output of each
execution of \texttt{svnadmin\ dump} is self-sufficient, the first
dumped revision is, by default, a full representation of every
directory, file, and property in that revision of the repository.

However, you can change this default behavior. If you add the
\texttt{-\/-incremental} option when you dump your repository,
\texttt{svnadmin} will compare the first dumped revision against the
previous revision in the repository---the same way it treats every other
revision that gets dumped. It will then output the first revision
exactly as it does the rest of the revisions in the dump
range---mentioning only the changes that occurred in that revision. The
benefit of this is that you can create several small dump files that can
be loaded in succession, instead of one large one, like so:

\begin{verbatim}
$ svnadmin dump myrepos -r 0:1000 > dumpfile1
$ svnadmin dump myrepos -r 1001:2000 --incremental > dumpfile2
$ svnadmin dump myrepos -r 2001:3000 --incremental > dumpfile3
\end{verbatim}

These dump files could be loaded into a new repository with the
following command sequence:

\begin{verbatim}
$ svnadmin load newrepos < dumpfile1
$ svnadmin load newrepos < dumpfile2
$ svnadmin load newrepos < dumpfile3
\end{verbatim}

Another neat trick you can perform with this \texttt{-\/-incremental}
option involves appending to an existing dump file a new range of dumped
revisions. For example, you might have a post-commit hook that simply
appends the repository dump of the single revision that triggered the
hook. Or you might have a script that runs nightly to append dump file
data for all the revisions that were added to the repository since the
last time the script ran. Used like this, \texttt{svnadmin\ dump} can be
one way to back up changes to your repository over time in case of a
system crash or some other catastrophic event.

The dump format can also be used to merge the contents of several
different repositories into a single repository. By using the
\texttt{-\/-parent-dir} option of \texttt{svnadmin\ load}, you can
specify a new virtual root directory for the load process. That means if
you have dump files for three repositories---say \texttt{calc-dumpfile},
\texttt{cal-dumpfile}, and \texttt{ss-dumpfile}---you can first create a
new repository to hold them all:

\begin{verbatim}
$ svnadmin create /var/svn/projects
$
\end{verbatim}

Then, make new directories in the repository that will encapsulate the
contents of each of the three previous repositories:

\begin{verbatim}
$ svn mkdir -m "Initial project roots" \
            file:///var/svn/projects/calc \
            file:///var/svn/projects/calendar \
            file:///var/svn/projects/spreadsheet
Committed revision 1.
$ 
\end{verbatim}

Lastly, load the individual dump files into their respective locations
in the new repository:

\begin{verbatim}
$ svnadmin load /var/svn/projects --parent-dir calc < calc-dumpfile
…
$ svnadmin load /var/svn/projects --parent-dir calendar < cal-dumpfile
…
$ svnadmin load /var/svn/projects --parent-dir spreadsheet < ss-dumpfile
…
$
\end{verbatim}

\subsubsection{Repository data migration using
svnrdump}\label{svn.reposadmin.maint.migrate.svnrdump}

In Subversion 1.7, \texttt{svnrdump} joined the set of stock Subversion
tools. It offers fairly specialized functionality, essentially as a
network-aware version of the \texttt{svnadmin\ dump} and
\texttt{svnadmin\ load} commands which we discuss in depth in
\hyperref[svn.reposadmin.maint.migrate.svnadmin]{Repository data
migration using svnadmin}. \texttt{svnrdump\ dump} will generate a dump
stream from a remote repository, spewing it to standard output;
\texttt{svnrdump\ load} will read a dump stream from standard input and
load it into a remote repository. Using \texttt{svnrdump}, you can
generate incremental dumps just as you might with
\texttt{svnadmin\ dump}. You can even dump a subtree of the
repository---something that \texttt{svnadmin\ dump} cannot do.

The primary difference is that instead of requiring direct access to the
repository, \texttt{svnrdump} operates remotely, using the very same
Repository Access (RA) protocols that the Subversion client does. As
such, you might need to provide authentication credentials. Also, your
remote interactions are subject to any authorization limitations
configured on the Subversion server.

\texttt{svnrdump\ dump} requires that the remote server be running
Subversion 1.4 or newer. It currently generates dump streams only of the
sort which are created when you pass the \texttt{-\/-deltas} option to
\texttt{svnadmin\ dump}. This isn\textquotesingle t interesting in the
typical use-cases, but might impact specific types of custom
transformations you might wish to apply to the resulting dump stream.

Because it modifies revision properties after committing new revisions,
\texttt{svnrdump\ load} requires that the target repository have
revision property changes enabled via the pre-revprop-change hook. See
\hyperref[svn.ref.reposhooks.pre-revprop-change]{pre-revprop-change} in
\hyperref[svn.ref.reposhooks]{Subversion Repository Hook Reference} for
details.

As you might expect, you can use \texttt{svnadmin} and \texttt{svnrdump}
in concert. You can, for example, use \texttt{svnrdump\ dump} to
generate a dump stream from a remote repository, and pipe the results
thereof through \texttt{svnadmin\ load} to copy all that repository
history into a local repository. Or you can do the reverse, copying
history from a local repository into a remote one.

By using \texttt{file://} URLs, \texttt{svnrdump} can also access local
repositories, but it will be doing so via Subversion\textquotesingle s
Repository Access (RA) abstraction layer---you\textquotesingle ll get
better performance out of \texttt{svnadmin} in such situations.

\subsection{Filtering Repository
History}\label{svn.reposadmin.maint.filtering}

Since Subversion stores your versioned history using, at the very least,
binary differencing algorithms and data compression (optionally in a
completely opaque database system), attempting manual tweaks is unwise
if not quite difficult, and at any rate strongly discouraged. And once
data has been stored in your repository, Subversion generally
doesn\textquotesingle t provide an easy way to remove that
data.\footnote{That\textquotesingle s rather the reason you use version
  control at all, right?} But inevitably, there will be times when you
would like to manipulate the history of your repository. You might need
to strip out all instances of a file that was accidentally added to the
repository (and shouldn\textquotesingle t be there for whatever
reason).\footnote{Conscious, cautious removal of certain bits of
  versioned data is actually supported by real use cases.
  That\textquotesingle s why an ``obliterate'' feature has been one of
  the most highly requested Subversion features, and one which the
  Subversion developers hope to soon provide.} Or, perhaps you have
multiple projects sharing a single repository, and you decide to split
them up into their own repositories. To accomplish tasks such as these,
administrators need a more manageable and malleable representation of
the data in their repositories---the Subversion repository dump format.

As we described earlier in
\hyperref[svn.reposadmin.maint.migrate]{Migrating Repository Data
Elsewhere}, the Subversion repository dump format is a human-readable
representation of the changes that you\textquotesingle ve made to your
versioned data over time. Use the \texttt{svnadmin\ dump} or
\texttt{svnrdump\ dump} command to generate the dump data, and
\texttt{svnadmin\ load} or \texttt{svnrdump\ load} to populate a new
repository with it. The great thing about the human-readability aspect
of the dump format is that, if you aren\textquotesingle t careless about
it, you can manually inspect and modify it. Of course, the downside is
that if you have three years\textquotesingle{} worth of repository
activity encapsulated in what is likely to be a very large dump file, it
could take you a long, long time to manually inspect and modify it.

That\textquotesingle s where \texttt{svndumpfilter} becomes useful. This
program acts as a path-based filter for repository dump streams. Simply
give it either a list of paths you wish to keep or a list of paths you
wish to not keep, and then pipe your repository dump data through this
filter. The result will be a modified stream of dump data that contains
only the versioned paths you (explicitly or implicitly) requested.

Let\textquotesingle s look at a realistic example of how you might use
this program. Earlier in this chapter (see
\hyperref[svn.reposadmin.projects.chooselayout]{Planning Your Repository
Organization}), we discussed the process of deciding how to choose a
layout for the data in your repositories---using one repository per
project or combining them, arranging stuff within your repository, and
so on. But sometimes after new revisions start flying in, you rethink
your layout and would like to make some changes. A common change is the
decision to move multiple projects that are sharing a single repository
into separate repositories for each project.

Our imaginary repository contains three projects: \texttt{calc},
\texttt{calendar}, and \texttt{spreadsheet}. They have been living
side-by-side in a layout like this:

\begin{verbatim}
/
   calc/
      trunk/
      branches/
      tags/
   calendar/
      trunk/
      branches/
      tags/
   spreadsheet/
      trunk/
      branches/
      tags/
\end{verbatim}

To get these three projects into their own repositories, we first dump
the whole repository:

\begin{verbatim}
$ svnadmin dump /var/svn/repos > repos-dumpfile
* Dumped revision 0.
* Dumped revision 1.
* Dumped revision 2.
* Dumped revision 3.
…
$
\end{verbatim}

Next, run that dump file through the filter, each time including only
one of our top-level directories. This results in three new dump files:

\begin{verbatim}
$ svndumpfilter include calc < repos-dumpfile > calc-dumpfile
…
$ svndumpfilter include calendar < repos-dumpfile > cal-dumpfile
…
$ svndumpfilter include spreadsheet < repos-dumpfile > ss-dumpfile
…
$
\end{verbatim}

At this point, you have to make a decision. Each of your dump files will
create a valid repository, but will preserve the paths exactly as they
were in the original repository. This means that even though you would
have a repository solely for your \texttt{calc} project, that repository
would still have a top-level directory named \texttt{calc}. If you want
your \texttt{trunk}, \texttt{tags}, and \texttt{branches} directories to
live in the root of your repository, you might wish to edit your dump
files, tweaking the \texttt{Node-path} and \texttt{Node-copyfrom-path}
headers so that they no longer have that first \texttt{calc/} path
component. Also, you\textquotesingle ll want to remove the section of
dump data that creates the \texttt{calc} directory. It will look
something like the following:

\begin{verbatim}
Node-path: calc
Node-action: add
Node-kind: dir
Content-length: 0
  
\end{verbatim}

If you do plan on manually editing the dump file to remove a top-level
directory, make sure your editor is not set to automatically convert
end-of-line characters to the native format (e.g.,
\texttt{\textbackslash{}r\textbackslash{}n} to
\texttt{\textbackslash{}n}), as the content will then not agree with the
metadata. This will render the dump file useless.

All that remains now is to create your three new repositories, and load
each dump file into the right repository, ignoring the UUID found in the
dump stream:

\begin{verbatim}
$ svnadmin create calc
$ svnadmin load --ignore-uuid calc < calc-dumpfile
<<< Started new transaction, based on original revision 1
     * adding path : Makefile ... done.
     * adding path : button.c ... done.
…
$ svnadmin create calendar
$ svnadmin load --ignore-uuid calendar < cal-dumpfile
<<< Started new transaction, based on original revision 1
     * adding path : Makefile ... done.
     * adding path : cal.c ... done.
…
$ svnadmin create spreadsheet
$ svnadmin load --ignore-uuid spreadsheet < ss-dumpfile
<<< Started new transaction, based on original revision 1
     * adding path : Makefile ... done.
     * adding path : ss.c ... done.
…
$
\end{verbatim}

Both of \texttt{svndumpfilter}\textquotesingle s subcommands accept
options for deciding how to deal with ``empty'' revisions. If a given
revision contains only changes to paths that were filtered out, that
now-empty revision could be considered uninteresting or even unwanted.
So to give the user control over what to do with those revisions,
\texttt{svndumpfilter} provides the following command-line options:

\begin{description}
\item[\texttt{-\/-drop-empty-revs}]
Do not generate empty revisions at all---just omit them.
\item[\texttt{-\/-renumber-revs}]
If empty revisions are dropped (using the \texttt{-\/-drop-empty-revs}
option), change the revision numbers of the remaining revisions so that
there are no gaps in the numeric sequence.
\item[\texttt{-\/-preserve-revprops}]
If empty revisions are not dropped, preserve the revision properties
(log message, author, date, custom properties, etc.) for those empty
revisions. Otherwise, empty revisions will contain only the original
datestamp, and a generated log message that indicates that this revision
was emptied by \texttt{svndumpfilter}.
\end{description}

While \texttt{svndumpfilter} can be very useful and a huge timesaver,
there are unfortunately a couple of gotchas. First, this utility is
overly sensitive to path semantics. Pay attention to whether paths in
your dump file are specified with or without leading slashes.
You\textquotesingle ll want to look at the \texttt{Node-path} and
\texttt{Node-copyfrom-path} headers.

\begin{verbatim}
…
Node-path: spreadsheet/Makefile
…
\end{verbatim}

If the paths have leading slashes, you should include leading slashes in
the paths you pass to \texttt{svndumpfilter\ include} and
\texttt{svndumpfilter\ exclude} (and if they don\textquotesingle t, you
shouldn\textquotesingle t). Further, if your dump file has an
inconsistent usage of leading slashes for some reason,\footnote{While
  \texttt{svnadmin\ dump} has a consistent leading slash policy (to not
  include them), other programs that generate dump data might not be so
  consistent.} you should probably normalize those paths so that they
all have, or all lack, leading slashes.

Also, copied paths can give you some trouble. Subversion supports copy
operations in the repository, where a new path is created by copying
some already existing path. It is possible that at some point in the
lifetime of your repository, you might have copied a file or directory
from some location that \texttt{svndumpfilter} is excluding, to a
location that it is including. To make the dump data self-sufficient,
\texttt{svndumpfilter} needs to still show the addition of the new
path---including the contents of any files created by the copy---and not
represent that addition as a copy from a source that
won\textquotesingle t exist in your filtered dump data stream. But
because the Subversion repository dump format shows only what was
changed in each revision, the contents of the copy source might not be
readily available. If you suspect that you have any copies of this sort
in your repository, you might want to rethink your set of
included/excluded paths, perhaps including the paths that served as
sources of your troublesome copy operations, too.

Finally, \texttt{svndumpfilter} takes path filtering quite literally. If
you are trying to copy the history of a project rooted at
\texttt{trunk/my-project} and move it into a repository of its own, you
would, of course, use the \texttt{svndumpfilter\ include} command to
keep all the changes in and under \texttt{trunk/my-project}. But the
resultant dump file makes no assumptions about the repository into which
you plan to load this data. Specifically, the dump data might begin with
the revision that added the \texttt{trunk/my-project} directory, but it
will \emph{not} contain directives that would create the \texttt{trunk}
directory itself (because \texttt{trunk} doesn\textquotesingle t match
the include filter). You\textquotesingle ll need to make sure that any
directories that the new dump stream expects to exist actually do exist
in the target repository before trying to load the stream into that
repository.

\subsection{Repository
Replication}\label{svn.reposadmin.maint.replication}

There are several scenarios in which it is quite handy to have a
Subversion repository whose version history is exactly the same as some
other repository\textquotesingle s. Perhaps the most obvious one is the
maintenance of a simple backup repository, used when the primary
repository has become inaccessible due to a hardware failure, network
outage, or other such annoyance. Other scenarios include deploying
mirror repositories to distribute heavy Subversion load across multiple
servers, use as a soft-upgrade mechanism, and so on.

Subversion provides a program for managing scenarios such as these.
\texttt{svnsync} works by essentially asking the Subversion server to
``replay'' revisions, one at a time. It then uses that revision
information to mimic a commit of the same to another repository. Neither
repository needs to be locally accessible to the machine on which
\texttt{svnsync} is running---its parameters are repository URLs, and it
does all its work through Subversion\textquotesingle s Repository Access
(RA) interfaces. All it requires is read access to the source repository
and read/write access to the destination repository.

When using \texttt{svnsync} against a remote source repository, the
Subversion server for that repository must be running Subversion version
1.4 or later.

\subsubsection{Replication with
svnsync}\label{svn.reposadmin.maint.replication.svnsync}

Assuming you already have a source repository that you\textquotesingle d
like to mirror, the next thing you need is a target repository that will
actually serve as that mirror. This target repository can use either of
the available filesystem data-store backends (see
\hyperref[svn.reposadmin.basics.backends]{Speaking of
Filesystems\ldots{}})---Subversion\textquotesingle s abstraction layers
ensure that such details don\textquotesingle t matter. But by default,
it must not yet have any version history in it. (We\textquotesingle ll
discuss an exception to this later in this section.)

The protocol that \texttt{svnsync} uses to communicate revision
information is highly sensitive to mismatches between the versioned
histories contained in the source and target repositories. For this
reason, while \texttt{svnsync} cannot \emph{demand} that the target
repository be read-only,\footnote{In fact, it can\textquotesingle t
  truly be read-only, or \texttt{svnsync} itself would have a tough time
  copying revision history into it.} allowing the revision history in
the target repository to change by any mechanism other than the
mirroring process is a recipe for disaster.

Do \emph{not} modify a mirror repository in such a way as to cause its
version history to deviate from that of the repository it mirrors. The
only commits and revision property modifications that ever occur on that
mirror repository should be those performed by the \texttt{svnsync}
tool.

Another requirement of the target repository is that the
\texttt{svnsync} process be allowed to modify revision properties.
Because \texttt{svnsync} works within the framework of that
repository\textquotesingle s hook system, the default state of the
repository (which is to disallow revision property changes; see
\hyperref[svn.ref.reposhooks.pre-revprop-change]{pre-revprop-change} in
\hyperref[svn.ref.reposhooks]{Subversion Repository Hook Reference}) is
insufficient. You\textquotesingle ll need to explicitly implement the
pre-revprop-change hook, and your script must allow \texttt{svnsync} to
set and change revision properties. With those provisions in place, you
are ready to start mirroring repository revisions.

It\textquotesingle s a good idea to implement authorization measures
that allow your repository replication process to perform its tasks
while preventing other users from modifying the contents of your mirror
repository at all.

Let\textquotesingle s walk through the use of \texttt{svnsync} in a
somewhat typical mirroring scenario. We\textquotesingle ll pepper this
discourse with practical recommendations, which you are free to
disregard if they aren\textquotesingle t required by or suitable for
your environment.

We will be mirroring the public Subversion repository which houses the
source code for this very book and exposing that mirror publicly on the
Internet, hosted on a different machine than the one on which the
original Subversion source code repository lives. This remote host has a
global configuration that permits anonymous users to read the contents
of repositories on the host, but requires users to authenticate to
modify those repositories. (Please forgive us for glossing over the
details of Subversion server configuration for the moment---those are
covered thoroughly in \hyperref[svn.serverconfig]{Server
Configuration}.) And for no other reason than that it makes for a more
interesting example, we\textquotesingle ll be driving the replication
process from a third machine---the one that we currently find ourselves
using.

First, we\textquotesingle ll create the repository which will be our
mirror. This and the next couple of steps do require shell access to the
machine on which the mirror repository will live. Once the repository is
all configured, though, we shouldn\textquotesingle t need to touch it
directly again.

\begin{verbatim}
$ ssh admin@svn.example.com "svnadmin create /var/svn/svn-mirror"
admin@svn.example.com's password: ********
$
\end{verbatim}

At this point, we have our repository, and due to our
server\textquotesingle s configuration, that repository is now ``live''
on the Internet. Now, because we don\textquotesingle t want anything
modifying the repository except our replication process, we need a way
to distinguish that process from other would-be committers. To do so, we
use a dedicated username for our process. Only commits and revision
property modifications performed by the special username
\texttt{syncuser} will be allowed.

We\textquotesingle ll use the repository\textquotesingle s hook system
both to allow the replication process to do what it needs to do and to
enforce that only it is doing those things. We accomplish this by
implementing two of the repository event hooks---pre-revprop-change and
start-commit. Our pre-revprop-change hook script is found in
\hyperref[svn.reposadmin.maint.replication.pre-revprop-change]{Mirror
repository\textquotesingle s pre-revprop-change hook script}, and
basically verifies that the user attempting the property changes is our
\texttt{syncuser} user. If so, the change is allowed; otherwise, it is
denied.

\protect\phantomsection\label{svn.reposadmin.maint.replication.pre-revprop-change}
Mirror repository\textquotesingle s pre-revprop-change hook script

\begin{verbatim}
#!/bin/sh 

USER="$3"

if [ "$USER" = "syncuser" ]; then exit 0; fi

echo "Only the syncuser user may change revision properties" >&2
exit 1
\end{verbatim}

That covers revision property changes. Now we need to ensure that only
the \texttt{syncuser} user is permitted to commit new revisions to the
repository. We do this using a start-commit hook script such as the one
in \hyperref[svn.reposadmin.maint.replication.start-commit]{Mirror
repository\textquotesingle s start-commit hook script}.

\protect\phantomsection\label{svn.reposadmin.maint.replication.start-commit}
Mirror repository\textquotesingle s start-commit hook script

\begin{verbatim}
#!/bin/sh 

USER="$2"

if [ "$USER" = "syncuser" ]; then exit 0; fi

echo "Only the syncuser user may commit new revisions" >&2
exit 1
\end{verbatim}

After installing our hook scripts and ensuring that they are executable
by the Subversion server, we\textquotesingle re finished with the setup
of the mirror repository. Now, we get to actually do the mirroring.

The first thing we need to do with \texttt{svnsync} is to register in
our target repository the fact that it will be a mirror of the source
repository. We do this using the \texttt{svnsync\ initialize}
subcommand. The URLs we provide point to the root directories of the
target and source repositories, respectively. In Subversion 1.4, this is
required---only full mirroring of repositories is permitted. Beginning
with Subversion 1.5, though, you can use \texttt{svnsync} to mirror only
some subtree of the repository, too.

\begin{verbatim}
$ svnsync help init
initialize (init): usage: svnsync initialize DEST_URL SOURCE_URL

Initialize a destination repository for synchronization from
another repository.
…
$ svnsync initialize http://svn.example.com/svn-mirror \
                     https://svn.code.sf.net/p/svnbook/source \
                     --sync-username syncuser --sync-password syncpass
Copied properties for revision 0 (svn:sync-* properties skipped).
NOTE: Normalized svn:* properties to LF line endings (1 rev-props, 0 node-props).
$
\end{verbatim}

Our target repository will now remember that it is a mirror of the
public Subversion source code repository. Notice that we provided a
username and password as arguments to \texttt{svnsync}---that was
required by the pre-revprop-change hook on our mirror repository.

In Subversion 1.4, the values given to
\texttt{svnsync}\textquotesingle s \texttt{-\/-username} and
\texttt{-\/-password} command-line options were used for authentication
against both the source and destination repositories. This caused
problems when a user\textquotesingle s credentials
weren\textquotesingle t exactly the same for both repositories,
especially when running in noninteractive mode (with the
\texttt{-\/-non-interactive} option). This was fixed in Subversion 1.5
with the introduction of two new pairs of options. Use
\texttt{-\/-source-username} and \texttt{-\/-source-password} to provide
authentication credentials for the source repository; use
\texttt{-\/-sync-username} and \texttt{-\/-sync-password} to provide
credentials for the destination repository. (The old
\texttt{-\/-username} and \texttt{-\/-password} options still exist for
compatibility, but we advise against using them.)

And now comes the fun part. With a single subcommand, we can tell
\texttt{svnsync} to copy all the as-yet-unmirrored revisions from the
source repository to the target.\footnote{Be forewarned that while it
  will take only a few seconds for the average reader to parse this
  paragraph and the sample output that follows it, the actual time
  required to complete such a mirroring operation is, shall we say,
  quite a bit longer.} The \texttt{svnsync\ synchronize} subcommand will
peek into the special revision properties previously stored on the
target repository and determine how much of the source repository has
been previously mirrored---in this case, the most recently mirrored
revision is r0. Then it will query the source repository and determine
what the latest revision in that repository is. Finally, it asks the
source repository\textquotesingle s server to start replaying all the
revisions between 0 and that latest revision. As \texttt{svnsync} gets
the resultant response from the source repository\textquotesingle s
server, it begins forwarding those revisions to the target
repository\textquotesingle s server as new commits.

\begin{verbatim}
$ svnsync help synchronize
synchronize (sync): usage: svnsync synchronize DEST_URL [SOURCE_URL]

Transfer all pending revisions to the destination from the source
with which it was initialized.
…
$ svnsync synchronize http://svn.example.com/svn-mirror \
                      https://svn.code.sf.net/p/svnbook/source
Committed revision 1.
Copied properties for revision 1.
Committed revision 2.
Copied properties for revision 2.
Transmitting file data .
Committed revision 3.
Copied properties for revision 3.
…
Transmitting file data .
Committed revision 4063.
Copied properties for revision 4063.
Transmitting file data .
Committed revision 4064.
Copied properties for revision 4064.
Transmitting file data ....
Committed revision 4065.
Copied properties for revision 4065.
$
\end{verbatim}

Of particular interest here is that for each mirrored revision, there is
first a commit of that revision to the target repository, and then
property changes follow. This two-phase replication is required because
the initial commit is performed by (and attributed to) the user
\texttt{syncuser} and is datestamped with the time as of that
revision\textquotesingle s creation. \texttt{svnsync} has to follow up
with an immediate series of property modifications that copy into the
target repository all the original revision properties found for that
revision in the source repository, which also has the effect of fixing
the author and datestamp of the revision to match that of the source
repository.

Also noteworthy is that \texttt{svnsync} performs careful bookkeeping
that allows it to be safely interrupted and restarted without ruining
the integrity of the mirrored data. If a network glitch occurs while
mirroring a repository, simply repeat the \texttt{svnsync\ synchronize}
command, and it will happily pick up right where it left off. In fact,
as new revisions appear in the source repository, this is exactly what
you do to keep your mirror up to date.

As part of its bookkeeping, \texttt{svnsync} records in the mirror
repository the URL with which the mirror was initialized. Because of
this, invocations of \texttt{svnsync} which follow the initialization
step do not \emph{require} that you provide the source URL on the
command line again. However, for security purposes, we recommend that
you continue to do so. Depending on how it is deployed, it may not be
safe for \texttt{svnsync} to trust the source URL which it retrieves
from the mirror repository, and from which it pulls versioned data.

svnsync Bookkeeping

\texttt{svnsync} needs to be able to set and modify revision properties
on the mirror repository because those properties are part of the data
it is tasked with mirroring. As those properties change in the source
repository, those changes need to be reflected in the mirror repository,
too. But \texttt{svnsync} also uses a set of custom revision
properties---stored in revision 0 of the mirror repository---for its own
internal bookkeeping. These properties contain information such as the
URL and UUID of the source repository, plus some additional
state-tracking information.

One of those pieces of state-tracking information is a flag that
essentially just means ``there\textquotesingle s a synchronization in
progress right now.'' This is used to prevent multiple \texttt{svnsync}
processes from colliding with each other while trying to mirror data to
the same destination repository. Now, generally you
won\textquotesingle t need to pay any attention whatsoever to \emph{any}
of these special properties (all of which begin with the prefix
\texttt{svn:sync-}). Occasionally, though, if a synchronization fails
unexpectedly, Subversion never has a chance to remove this particular
state flag. This causes all future synchronization attempts to fail
because it appears that a synchronization is still in progress when, in
fact, none is. Fortunately, recovering from this situation is easy to
do. In Subversion 1.7, you can use the newly introduced
\texttt{-\/-steal-lock} option with \texttt{svnsync}\textquotesingle s
commands. In previous Subversion versions, you need only to remove the
\texttt{svn:sync-lock} property which serves as this flag from revision
0 of the mirror repository:

\begin{verbatim}
$ svn propdel --revprop -r0 svn:sync-lock http://svn.example.com/svn-mirror
property 'svn:sync-lock' deleted from repository revision 0
$
\end{verbatim}

Also, \texttt{svnsync} stores the source repository URL provided at
mirror initialization time in a bookkeeping property on the mirror
repository. Future synchronization operations against that mirror which
omit the source URL at the command line will consult the special
\texttt{svn:sync-from-url} property stored on the mirror itself to know
where to synchronize from. This value is used literally by the
synchronization process, though. Be wary of using non-fully-qualified
domain names (such as referring to \texttt{svnbook.red-bean.com} as
simply \texttt{svnbook} because that happens to work when you are
connected directly to the \texttt{red-bean.com} network), domain names
which don\textquotesingle t resolve or resolve differently depending on
where you happen to be operating from, or IP addresses (which can change
over time). But here again, if you need an existing mirror to start
referring to a different URL for the same source repository, you can
change the bookkeeping property which houses that information. Users of
Subversion 1.7 or better can use
\texttt{svnsync\ init\ -\/-allow-non-empty} to reinitialize their
mirrors with new source URL:

\begin{verbatim}
$ svnsync initialize --allow-non-empty http://svn.example.com/svn-mirror \
                                       NEW-SOURCE-URL
Copied properties for revision 4065.
$
\end{verbatim}

If you are running an older version of Subversion,
you\textquotesingle ll need to manually tweak the
\texttt{svn:sync-from-url} bookkeeping property:

\begin{verbatim}
$ svn propset --revprop -r0 svn:sync-from-url NEW-SOURCE-URL \
      http://svn.example.com/svn-mirror
property 'svn:sync-from-url' set on repository revision 0
$
\end{verbatim}

Another interesting thing about these special bookkeeping properties is
that \texttt{svnsync} will not attempt to mirror any of those properties
when they are found in the source repository. The reason is probably
obvious, but basically boils down to \texttt{svnsync} not being able to
distinguish the special properties it has merely copied from the source
repository from those it needs to consult and maintain for its own
bookkeeping needs. This situation could occur if, for example, you were
maintaining a mirror of a mirror of a third repository. When
\texttt{svnsync} sees its own special properties in revision 0 of the
source repository, it simply ignores them.

An \texttt{svnsync\ info} subcommand was added in Subversion 1.6 to
easily display the special bookkeeping properties in the destination
repository.

\begin{verbatim}
$ svnsync help info
info: usage: svnsync info DEST_URL

Print information about the synchronization destination repository
located at DEST_URL.
…
$ svnsync info http://svn.example.com/svn-mirror
Source URL: https://svn.code.sf.net/p/svnbook/source
Source Repository UUID: 931749d0-5854-0410-9456-f14be4d6b398
Last Merged Revision: 4065
$
\end{verbatim}

There is, however, one bit of inelegance in the process. Because
Subversion revision properties can be changed at any time throughout the
lifetime of the repository, and because they don\textquotesingle t leave
an audit trail that indicates when they were changed, replication
processes have to pay special attention to them. If
you\textquotesingle ve already mirrored the first 15 revisions of a
repository and someone then changes a revision property on revision 12,
\texttt{svnsync} won\textquotesingle t know to go back and patch up its
copy of revision 12. You\textquotesingle ll need to tell it to do so
manually by using (or with some additional tooling around) the
\texttt{svnsync\ copy-revprops} subcommand, which simply rereplicates
all the revision properties for a particular revision or range thereof.

\begin{verbatim}
$ svnsync help copy-revprops
copy-revprops: usage:

    1. svnsync copy-revprops DEST_URL [SOURCE_URL]
    2. svnsync copy-revprops DEST_URL REV[:REV2]

…
$ svnsync copy-revprops http://svn.example.com/svn-mirror 12
Copied properties for revision 12.
$
\end{verbatim}

That\textquotesingle s repository replication via \texttt{svnsync} in a
nutshell. You\textquotesingle ll likely want some automation around such
a process. For example, while our example was a pull-and-push setup, you
might wish to have your primary repository push changes to one or more
blessed mirrors as part of its post-commit and post-revprop-change hook
implementations. This would enable the mirror to be up to date in as
near to real time as is likely possible.

\subsubsection{Partial replication with
svnsync}\label{svn.reposadmin.maint.replication.svnsync-partial}

\texttt{svnsync} isn\textquotesingle t limited to full copies of
everything which lives in a repository. It can handle various shades of
partial replication, too. For example, while it isn\textquotesingle t
very commonplace to do so, \texttt{svnsync} does gracefully mirror
repositories in which the user as whom it authenticates has only partial
read access. It simply copies only the bits of the repository that it is
permitted to see. Obviously, such a mirror is not useful as a backup
solution.

As of Subversion 1.5, \texttt{svnsync} also has the ability to mirror a
subset of a repository rather than the whole thing. The process of
setting up and maintaining such a mirror is exactly the same as when
mirroring a whole repository, except that instead of specifying the
source repository\textquotesingle s root URL when running
\texttt{svnsync\ init}, you specify the URL of some subdirectory within
that repository. Synchronization to that mirror will now copy only the
bits that changed under that source repository subdirectory. There are
some limitations to this support, though. First, you
can\textquotesingle t mirror multiple disjoint subdirectories of the
source repository into a single mirror
repository---you\textquotesingle d need to instead mirror some parent
directory that is common to both. Second, the filtering logic is
entirely path-based, so if the subdirectory you are mirroring was
renamed at some point in the past, your mirror would contain only the
revisions since the directory appeared at the URL you specified. And
likewise, if the source subdirectory is renamed in the future, your
synchronization processes will stop mirroring data at the point that the
source URL you specified is no longer valid.

\subsubsection{A quick trick for mirror
creation}\label{svn.reposadmin.maint.replication.svnsync-init-nonempty}

We mentioned previously the cost of setting up an initial mirror of an
existing repository. For many folks, the sheer cost of transmitting
thousands---or millions---of revisions of history to a new mirror
repository via \texttt{svnsync} is a show-stopper. Fortunately,
Subversion 1.7 provides a workaround by way of a new
\texttt{-\/-allow-non-empty} option to \texttt{svnsync\ initialize}.
This option allows you to initialize one repository as a mirror of
another while bypassing the verification that the to-be-initialized
mirror has no version history present in it. Per our previous warnings
about the sensitivity of this whole replication process, you should
rightly discern that this is an option to be used only with great
caution. But it\textquotesingle s wonderfully handy when you have
administrative access to the source repository, where you can simply
make a physical copy of the repository and then initialize that copy as
a new mirror:

\begin{verbatim}
$ svnadmin hotcopy /path/to/repos /path/to/mirror-repos
$ ### create /path/to/mirror-repos/hooks/pre-revprop-change
$ svnsync initialize file:///path/to/mirror-repos \
                     file:///path/to/repos
svnsync: E000022: Destination repository already contains revision history; co
nsider using --allow-non-empty if the repository's revisions are known to mirr
or their respective revisions in the source repository
$ svnsync initialize --allow-non-empty file:///path/to/mirror-repos \
                                       file:///path/to/repos
Copied properties for revision 32042.
$
\end{verbatim}

Admins who are running a version of Subversion prior to 1.7 (and thus do
not have access to \texttt{svnsync\ initialize}\textquotesingle s
\texttt{-\/-allow-non-empty} feature) can accomplish effectively the
same thing that that feature does through \emph{careful} manipulation of
the r0 revision properties on the copy of the repository which is slated
to become a mirror of the original. Use \texttt{svnadmin\ setrevprop} to
create the same bookkeeping properties that \texttt{svnsync} would have
created there.

\subsubsection{Replication
wrap-up}\label{svn.reposadmin.maint.replication.wrapup}

We\textquotesingle ve discussed a couple of ways to replicate revision
history from one repository to another. So let\textquotesingle s look
now at the user end of these operations. How does replication and the
various situations which call for it affect Subversion clients?

As far as user interaction with repositories and mirrors goes, it
\emph{is} possible to have a single working copy that interacts with
both, but you\textquotesingle ll have to jump through some hoops to make
it happen. First, you need to ensure that both the primary and mirror
repositories have the same repository UUID (which is not the case by
default). See \hyperref[svn.reposadmin.maint.uuids]{Managing Repository
UUIDs} later in this chapter for more about this.

Once the two repositories have the same UUID, you can use
\texttt{svn\ relocate} to point your working copy to whichever of the
repositories you wish to operate against, a process that is described in
\hyperref[svn.ref.svn.c.relocate]{svn relocate} in
\hyperref[svn.ref.svn]{svn Reference---Subversion Command-Line Client}.
There is a possible danger here, though, in that if the primary and
mirror repositories aren\textquotesingle t in close synchronization, a
working copy up to date with, and pointing to, the primary repository
will, if relocated to point to an out-of-date mirror, become confused
about the apparent sudden loss of revisions it fully expects to be
present, and it will throw errors to that effect. If this occurs, you
can relocate your working copy back to the primary repository and then
either wait until the mirror repository is up to date, or backdate your
working copy to a revision you know is present in the sync repository,
and then retry the relocation.

Finally, be aware that the revision-based replication provided by
\texttt{svnsync} is only that---replication of revisions. Only the kinds
of information carried by the Subversion repository dump file format are
available for replication. As such, tools such as \texttt{svnsync} (and
\texttt{svnrdump}, which we discuss in
\hyperref[svn.reposadmin.maint.migrate.svnrdump]{Repository data
migration using svnrdump}) are limited in ways similar to that of the
repository dump stream. They do not include in their replicated
information such things as the hook implementations, repository or
server configuration data, uncommitted transactions, or information
about user locks on repository paths.

\subsection{Repository Backup}\label{svn.reposadmin.maint.backup}

Despite numerous advances in technology since the birth of the modern
computer, one thing unfortunately rings true with crystalline
clarity---sometimes things go very, very awry. Power outages, network
connectivity dropouts, corrupt RAM, and crashed hard drives are but a
taste of the evil that Fate is poised to unleash on even the most
conscientious administrator. And so we arrive at a very important
topic---how to make backup copies of your repository data.

There are two types of backup methods available for Subversion
repository administrators---full and incremental. A full backup of the
repository involves squirreling away in one sweeping action all the
information required to fully reconstruct that repository in the event
of a catastrophe. Usually, it means, quite literally, the duplication of
the entire repository directory (which includes the FSFS
environment). Incremental backups are lesser things: backups of
only the portion of the repository data that has changed since the
previous backup.

As far as full backups go, the naïve approach might seem like a sane
one, but unless you temporarily disable all other access to your
repository, simply doing a recursive directory copy runs the risk of
generating a faulty backup. Subversion\textquotesingle s FSFS data must
be copied in a particular order to guarantee a valid backup copy. But
you don\textquotesingle t have to implement that algorithm yourself,
because the Subversion development team has already done so. The \texttt{svnadmin\ hotcopy} command takes care of
the minutiae involved in making a hot backup of your repository. And its
invocation is as trivial as the Unix \texttt{cp} or Windows
\texttt{copy} operations:

\begin{verbatim}
$ svnadmin hotcopy /var/svn/repos /var/svn/repos-backup
\end{verbatim}

The resultant backup is a fully functional Subversion repository, able
to be dropped in as a replacement for your live repository should
something go horribly wrong.

Additional tooling around this command is available, too. The
\texttt{tools/backup/} directory of the Subversion source distribution
holds the \texttt{hot-backup.py} script. This script adds a bit of
backup management atop \texttt{svnadmin\ hotcopy}, allowing you to keep
only the most recent configured number of backups of each repository. It
will automatically manage the names of the backed-up repository
directories to avoid collisions with previous backups and will ``rotate
off'' older backups, deleting them so that only the most recent ones
remain. Even if you also have an incremental backup, you might want to
run this program on a regular basis. For example, you might consider
using \texttt{hot-backup.py} from a program scheduler (such as
\texttt{cron} on Unix systems), which can cause it to run nightly (or at
whatever granularity of time you deem safe).

Some administrators use a different backup mechanism built around
generating and storing repository dump data. We described in
\hyperref[svn.reposadmin.maint.migrate]{Migrating Repository Data
Elsewhere} how to use \texttt{svnadmin\ dump} with the
\texttt{-\/-incremental} option to perform an incremental backup of a
given revision or range of revisions. And of course, you can achieve a
full backup variation of this by omitting the \texttt{-\/-incremental}
option to that command. There is some value in these methods, in that
the format of your backed-up information is
flexible---it\textquotesingle s not tied to a particular platform,
versioned filesystem type, or release of Subversion or the libraries it
uses. But that flexibility comes at a cost, namely that restoring that
data can take a long time---longer with each new revision committed to
your repository. Also, as is the case with so many of the various backup
methods, revision property changes that are made to already backed-up
revisions won\textquotesingle t get picked up by a nonoverlapping,
incremental dump generation. For these reasons, we recommend against
relying solely on dump-based backup approaches.

Beginning with Subversion 1.8, \texttt{svnadmin\ hotcopy} accepts
\texttt{-\/-incremental} option and supports incremental hotcopy mode
for FSFS repositories. In incremental hotcopy mode, revision data which
has already been copied from the source to the destination repository
will not be copied again. When \texttt{-\/-incremental} option is used
with \texttt{svnadmin\ hotcopy}, Subversion will only copy new
revisions, and revisions which have changed in size or had their
modification time stamp changed since the previous hotcopy operation.
Moreover, unlike with \texttt{svnsync} or
\texttt{svnadmin\ dump\ -\/-incremental}, performance of
\texttt{svnadmin\ hotcopy\ -\/-incremental} is only limited to disk I/O.
Therefore, incremental hotcopy can be a huge time saver when making a
backup of a large repository.

As you can see, each of the various backup types and methods has its
advantages and disadvantages. The easiest is by far the full hot backup,
which will always result in a perfect working replica of your
repository. Should something bad happen to your live repository, you can
restore from the backup with a simple recursive directory copy.
Unfortunately, if you are maintaining multiple backups of your
repository, these full copies will each eat up just as much disk space
as your live repository. Incremental backups, by contrast, tend to be
quicker to generate and smaller to store. But the restoration process
can be a pain, often involving applying multiple incremental backups.
And other methods have their own peculiarities. Administrators need to
find the balance between the cost of making the backup and the cost of
restoring it.

The \texttt{svnsync} program (see
\hyperref[svn.reposadmin.maint.replication]{Repository Replication})
actually provides a rather handy middle-ground approach. If you are
regularly synchronizing a read-only mirror with your main repository, in
a pinch your read-only mirror is probably a good candidate for replacing
that main repository if it falls over. The primary disadvantage of this
method is that only the versioned repository data gets
synchronized---repository configuration files, user-specified repository
path locks, and other items that might live in the physical repository
directory but not \emph{inside} the repository\textquotesingle s virtual
versioned filesystem are not handled by \texttt{svnsync}.

In any backup scenario, repository administrators need to be aware of
how modifications to unversioned revision properties affect their
backups. Since these changes do not themselves generate new revisions,
they will not trigger post-commit hooks, and may not even trigger the
pre-revprop-change and post-revprop-change hooks.\footnote{\texttt{svnadmin\ setlog}
  can be called in a way that bypasses the hook interface altogether.}
And since you can change revision properties without respect to
chronological order---you can change any revision\textquotesingle s
properties at any time---an incremental backup of the latest few
revisions might not catch a property modification to a revision that was
included as part of a previous backup.

Generally speaking, only the truly paranoid would need to back up their
entire repository, say, every time a commit occurred. However, assuming
that a given repository has some other redundancy mechanism in place
with relatively fine granularity (such as per-commit emails or
incremental dumps), a hot backup of the database might be something that
a repository administrator would want to include as part of a
system-wide nightly backup. It\textquotesingle s your data---protect it
as much as you\textquotesingle d like.

Often, the best approach to repository backups is a diversified one that
leverages combinations of the methods described here. The Subversion
developers, for example, back up the Subversion source code repository
nightly using \texttt{hot-backup.py} and an off-site \texttt{rsync} of
those full backups; keep multiple archives of all the commit and
property change notification emails; and have repository mirrors
maintained by various volunteers using \texttt{svnsync}. Your solution
might be similar, but should be catered to your needs and that delicate
balance of convenience with paranoia. And whatever you do, validate your
backups from time to time---what good is a spare tire that has a hole in
it? While all of this might not save your hardware from the iron fist of
Fate,\footnote{You know---the collective term for all of her ``fickle
  fingers.''} it should certainly help you recover from those trying
times.

\subsection{Managing Repository UUIDs}\label{svn.reposadmin.maint.uuids}

Subversion repositories have a universally unique identifier (UUID)
associated with them. This is used by Subversion clients to verify the
identity of a repository when other forms of verification
aren\textquotesingle t good enough (such as checking the repository URL,
which can change over time). Most Subversion repository administrators
rarely, if ever, need to think about repository UUIDs as anything more
than a trivial implementation detail of Subversion. Sometimes, however,
there is cause for attention to this detail.

As a general rule, you want the UUIDs of your live repositories to be
unique. That is, after all, the point of having UUIDs. But there are
times when you want the repository UUIDs of two repositories to be
exactly the same. For example, if you make a copy of a repository for
backup purposes, you want the backup to be a perfect replica of the
original so that, in the event that you have to restore that backup and
replace the live repository, users don\textquotesingle t suddenly see
what looks like a different repository. When dumping and loading
repository history (as described earlier in
\hyperref[svn.reposadmin.maint.migrate]{Migrating Repository Data
Elsewhere}), you get to decide whether to apply the UUID encapsulated in
the data dump stream to the repository in which you are loading the
data. The particular circumstance will dictate the correct behavior.

There are a couple of ways to set (or reset) a
repository\textquotesingle s UUID, should you need to. As of Subversion
1.5, this is as simple as using the \texttt{svnadmin\ setuuid} command.
If you provide this subcommand with an explicit UUID, it will validate
that the UUID is well-formed and then set the repository UUID to that
value. If you omit the UUID, a brand-new UUID will be generated for your
repository.

\begin{verbatim}
$ svnlook uuid /var/svn/repos
cf2b9d22-acb5-11dc-bc8c-05e83ce5dbec
$ svnadmin setuuid /var/svn/repos   # generate a new UUID
$ svnlook uuid /var/svn/repos
3c3c38fe-acc0-11dc-acbc-1b37ff1c8e7c
$ svnadmin setuuid /var/svn/repos \
           cf2b9d22-acb5-11dc-bc8c-05e83ce5dbec  # restore the old UUID
$ svnlook uuid /var/svn/repos
cf2b9d22-acb5-11dc-bc8c-05e83ce5dbec
$
\end{verbatim}

For folks using versions of Subversion earlier than 1.5, these tasks are
a little more complicated. You can explicitly set a
repository\textquotesingle s UUID by piping a repository dump file stub
that carries the new UUID specification through
\texttt{svnadmin\ load\ -\/-force-uuid\ REPOS-PATH}.

\begin{verbatim}
$ svnadmin load --force-uuid /var/svn/repos <<EOF
SVN-fs-dump-format-version: 2

UUID: cf2b9d22-acb5-11dc-bc8c-05e83ce5dbec
EOF
$ svnlook uuid /var/svn/repos
cf2b9d22-acb5-11dc-bc8c-05e83ce5dbec
$
\end{verbatim}

Having older versions of Subversion generate a brand-new UUID is not
quite as simple to do, though. Your best bet here is to find some other
way to generate a UUID, and then explicitly set the
repository\textquotesingle s UUID to that value.

\section{Moving and Removing
Repositories}\label{svn.reposadmin.maint.moving-and-removing}

Subversion repository data is wholly contained within the repository
directory. As such, you can move a Subversion repository to some other
location on disk, rename a repository, copy a repository, or delete a
repository altogether using the tools provided by your operating system
for manipulating directories---\texttt{mv}, \texttt{cp\ -a}, and
\texttt{rm\ -r} on Unix platforms; \texttt{copy}, \texttt{move}, and
\texttt{rmdir\ /s\ /q} on Windows; vast numbers of mouse and menu
gyrations in various graphical file explorer applications, and so on.

Of course, there\textquotesingle s often still more to be done when
trying to cleanly affect changes such as this. For example, you might
need to update your Subversion server configuration to point to the new
location of a relocated repository or to remove configuration bits for a
now-deleted repository. If you have automated processes that publish
information from or about your repositories, they may need to be
updated. Hook scripts might need to be reconfigured. Users may need to
be notified. The list can go on indefinitely, or at least to the extent
that you\textquotesingle ve built processes and procedures around your
Subversion repository.

In the case of a copied repository, you should also consider the fact
that Subversion uses repository UUIDs to distinguish repositories. If
you copy a Subversion repository using a typical shell recursive copy
command, you\textquotesingle ll wind up with two repositories that are
identical in every way---including their UUIDs. In some circumstances,
this might be desirable. But in the instances where it is not,
you\textquotesingle ll need to generate a new UUID for one of these
identical repositories. See
\hyperref[svn.reposadmin.maint.uuids]{Managing Repository UUIDs} for
more about managing repository UUIDs.

\section{Summary}\label{svn.reposadmin.summary}

By now you should have a basic understanding of how to create,
configure, and maintain Subversion repositories. We introduced you to
the various tools that will assist you with this task. Throughout the
chapter, we noted common administration pitfalls and offered suggestions
for avoiding them.

All that remains is for you to decide what exciting data to store in
your repository, and finally, how to make it available over a network.
The next chapter is all about networking.

